\documentclass{book}
%\documentclass{book}
%If you're going to print this, perhaps get ride of oneside so that \parts{} will have their own page!
\title{
	\begin{huge}
		\bf{Preliminary Chapter}
	\end{huge}
	}
\author{Nathanael Chwojko-Srawley}

%\date{} % If you want to specify a date

%  ╔══════════════════════════════════════════════════════════╗
%  ║                         packages                         ║
%  ╚══════════════════════════════════════════════════════════╝

\usepackage[a4paper, total={6in, 8in}]{geometry}
\usepackage[answerdelayed]{exercise}
\usepackage[font=itshape]{quoting}
\usepackage[most]{tcolorbox}         % Makes nice boxes for thms, cor, prop, or anything else.
\usepackage{amsmath, amssymb, amsthm} % Used to write better mathematical expressions (ex. \mathbb{R})
\usepackage{array}                   % \begin{table}{>{\em}c c} -> 1st column italic (bfseries for bold)
\usepackage{bbm}                     % for mathbbm{1}
\usepackage{blindtext}
\usepackage{commutative-diagrams}
\usepackage{dsfont}
\usepackage{dynkin-diagrams}         % for Dynkin diagrams
\usepackage{enumitem}
\usepackage{etoolbox}
\usepackage{euscript}
\usepackage{extarrows}               % Used to compile any UTF-8 Character (and supports Unicode)
\usepackage{fancyhdr}
\usepackage{float}                   % package with ``H'' for figures and tables, ex. Images, tables, etc
\usepackage{graphicx}                % Let's you use images in LaTeX; ex the \includegraphics command.
\usepackage{hhline}
\usepackage{hyperref}
\usepackage{cleveref}
\usepackage{inputenc}                % Used to compile any UTF-8 Character (and supports Unicode)
\usepackage{listings}                % Create nice colour code blocks (customized in preamble)
\usepackage{longtable}               % create tables that extend past one page
\usepackage{makeidx}
\usepackage{mathrsfs}
\usepackage{mathtools}
% \usepackage{microtype}               % makes nice text. Fixes some under-/over- filled box problems.
\usepackage{multicol}                % For multiple columns
\usepackage{pgfplots}
\usepackage{scalerel}
\usepackage{setspace}                % More flexibility on spacing in document.
\usepackage{stackengine}
\usepackage{stackrel}                % for left-right arrows, staking symbol above and bellow
\usepackage{thmtools}                % Customize theorem environment (in preamble).
\usepackage{tikz-cd}                 % for commutative diagrams https://tex.stackexchange.com/questions/419221/how-to-do-the-pushout-with-universal-property
\usepackage{tikz}                    % for commutative diagrams https://tex.stackexchange.com/questions/419221/how-to-do-the-pushout-with-universal-property
\usepackage{titlesec}
\usepackage{wasysym}                 % emoji's!
\usepackage{wrapfig}                 % To wrap text around figure
\usepackage{xcolor}                  % Colour text.

% \usepackage{comment}

% \usepackage{CJKutf8} % for some chinese

% \usepackage{dutchcal} % for lower case mathcal
% \pdfsuppresswarningpagegroup=1

% ╔═════════════════════════════════════════════════════╗
% ║           for figures via inkspace                  ║
% ╚═════════════════════════════════════════════════════╝

% to import figures via: https://github.com/gillescastel/inkscape-figures
\usepackage{import}
\usepackage{pdfpages}
\usepackage{transparent}


% This command is the ONLY command imported via the packages snippet, think of it as a tiny package written out
\newcommand{\incfig}[2][1]{%
    \def\svgwidth{#1\columnwidth}
    \import{./figures/}{#2.pdf_tex}
}

% ╔═══════════════════════════════════════════════════════════╗
% ║                 bibliography and citing                   ║
% ╚═══════════════════════════════════════════════════════════╝

\usepackage{csquotes}
% \usepackage[backend=biber,citestyle=alphabetic,bibstyle=authortitle]{biblatex}
\usepackage[backend=biber]{biblatex}

\addbibresource{bibliography.bib}

% ╔═══════════════════════════════════════════════════════════╗
% ║                    colours (colors)                       ║
% ╚═══════════════════════════════════════════════════════════╝

\definecolor{dkgreen}{rgb}{0,0.6,0}
\definecolor{gray}{rgb}{0.5,0.5,0.5}
\definecolor{mauve}{rgb}{0.58,0,0.82}
\definecolor{lightyellow}{rgb}{1,1,0.6}
\definecolor{reddish}{HTML}{A01A3A}

%  ╔══════════════════════════════════════════════════════════╗
%  ║               title/heading/footer format                ║
%  ╚══════════════════════════════════════════════════════════╝

\pagestyle{fancy}
\setlength\headheight{20pt}
\renewcommand{\sectionmark}[1]{\markright{\thesection.\ #1}}
\renewcommand{\chaptermark}[1]{\markboth{#1}{}}
\fancyfoot[C,CO]{\textcolor{gray}{Nathanael Chwojko-Srawley}}
\fancyfoot[R,RO]{\thepage}{}

\titleformat
{\chapter} % command
[display] % shape
{\bfseries\Huge\itshape} % format
{\thechapter} % label
{0.5ex} % sep
{
    \rule{\textwidth}{1pt}
    \vspace{1ex}
    \centering
}
[
\vspace{-0.5ex}
\rule{\textwidth}{0.3pt}
]

\setlength{\parindent}{0pt} % don't like indentation infront of paragraphs
\setlength{\parskip}{1ex}   %so that there is still space between paragarphs


%  ╔══════════════════════════════════════════════════════════╗
%  ║                       Shortcuts                          ║
%  ╚══════════════════════════════════════════════════════════╝

%\ExerciseLevelInToc

% I like original mathcal better, but need lowercase.
\DeclareFontFamily{OT1}{pzc}{}
\DeclareFontShape{OT1}{pzc}{m}{it}{<-> s * [1.10] pzcmi7t}{}
\DeclareMathAlphabet{\mathpzc}{OT1}{pzc}{m}{it}


% mathbb
\newcommand{\A}{{\mathbb{A}}}
\newcommand{\C}{{\mathbb{C}}}
\newcommand{\F}{{\mathbb{F}}}
\newcommand{\N}{{\mathbb{N}}}
\newcommand{\Q}{{\mathbb{Q}}}
\newcommand{\R}{{\mathbb{R}}}
\newcommand{\V}{{\mathbb{V}}}
\newcommand{\Z}{{\mathbb{Z}}}
\newcommand{\BI}{{\mathbb{I}}}
\renewcommand{\H}{{\mathbb{H}}}
\renewcommand{\P}{{\mathbb{P}}}


% mathcal
% \newcommand{\co}{{\mathfrak{o}}} %mathcal{o} looks like wreath product, so I changed it to mathfrak
\newcommand{\co}{{\mathpzc{o}}}
\newcommand{\CA}{{\mathcal{A}}}
\newcommand{\CB}{{\mathcal{B}}}
\newcommand{\CC}{{\mathcal{C}}}
\newcommand{\CD}{{\mathcal{D}}}
\newcommand{\CF}{{\mathcal{F}}}
\newcommand{\CG}{{\mathcal{G}}}
\newcommand{\CH}{{\mathcal{H}}}
\newcommand{\CI}{{\mathcal{I}}}
\newcommand{\CK}{{\mathcal{K}}}
\newcommand{\CL}{{\mathcal{L}}}
\newcommand{\CM}{{\mathcal{M}}}
\newcommand{\CN}{{\mathcal{N}}}
\newcommand{\CO}{{\mathcal{O}}}
\newcommand{\CQ}{{\mathcal{Q}}}
\newcommand{\CR}{{\mathcal{R}}}
\newcommand{\CS}{{\mathcal{S}}}
\newcommand{\CT}{{\mathcal{T}}}
\newcommand{\CV}{{\mathcal{V}}}
\newcommand{\CZ}{{\mathcal{Z}}}
% EXCPETION:
\newcommand{\CP}{{\mathcal{P}}}
% \newcommand{\CP}{\mathbb{C}\mathrm{P}}
\newcommand{\RP}{\mathbb{R}\mathrm{P}}


%Euscript
\newcommand{\EB}{\EuScript{B}}
\newcommand{\EC}{{\EuScript{C}}}
\newcommand{\EF}{{\EuScript{F}}}
\newcommand{\EL}{{\EuScript{L}}}
\newcommand{\EO}{{\EuScript{O}}}


%mathfrak (lowercase)
\newcommand{\fa}{{\mathfrak{a}}}
\newcommand{\fb}{{\mathfrak{b}}}
\newcommand{\fc}{{\mathfrak{c}}}
\newcommand{\fd}{{\mathfrak{d}}}
\newcommand{\fm}{{\mathfrak{m}}}
\newcommand{\fn}{{\mathfrak{n}}}
\newcommand{\fo}{{\mathfrak{o}}}
\newcommand{\fp}{{\mathfrak{p}}}
\newcommand{\fq}{{\mathfrak{q}}}


%mathfrak (upper case)
\newcommand{\FA}{{\mathfrak{A}}}
\newcommand{\FB}{{\mathfrak{B}}}
\newcommand{\FC}{{\mathfrak{C}}}
\newcommand{\FD}{{\mathfrak{D}}}
\newcommand{\FK}{{\mathfrak{K}}}
\newcommand{\FM}{{\mathfrak{M}}}
\newcommand{\FN}{{\mathfrak{N}}}
\newcommand{\FO}{{\mathfrak{O}}}
\newcommand{\FP}{{\mathfrak{P}}}


%mathscr
\newcommand{\ff}{{\mathscr{F}}}
\newcommand{\SA}{\mathscr{A}}
\newcommand{\SC}{\mathscr{C}}
\newcommand{\SF}{\mathscr{F}}
\newcommand{\SG}{\mathscr{G}}
\newcommand{\SH}{\mathscr{H}}
\newcommand{\SI}{\mathscr{I}}
\newcommand{\SK}{\mathscr{K}}
\newcommand{\SN}{\mathscr{N}}
\newcommand{\SQ}{\mathscr{Q}}
\newcommand{\SR}{\mathscr{R}}
\newcommand{\SV}{\mathscr{V}}
\newcommand{\SZ}{\mathscr{Z}}


% operatorname -- 2 letters (lowercase and upper case)
\newcommand{\ab}{{\operatorname{ab}}}
\newcommand{\ad}{{\operatorname{ad}}}
\newcommand{\ch}{{\operatorname{ch}}}
\newcommand{\ev}{{\operatorname{ev}}}
\newcommand{\id}{{\operatorname{id}}}
\newcommand{\im}{{\operatorname{im}}}
\newcommand{\nm}{{\operatorname{Nm}}}
\newcommand{\op}{{\operatorname{op}}}
\newcommand{\rk}{{\operatorname{rk}}}
\newcommand{\sh}{{\operatorname{sh}}}
\newcommand{\tr}{{\operatorname{tr}}}

\newcommand{\Ab}{{\operatorname{Ab}}}
\newcommand{\Br}{{\operatorname{Br}}}
\newcommand{\Ch}{{\operatorname{Ch}}}
\newcommand{\Cl}{{\operatorname{Cl}}}
\newcommand{\GL}{{\operatorname{GL}}}
\newcommand{\Nm}{{\operatorname{Nm}}}
\newcommand{\SL}{{\operatorname{SL}}}

\renewcommand{\Re}{{\operatorname{Re}}}
\renewcommand{\Im}{{\operatorname{Im}}}


%categories
\newcommand{\Grp}{{\textbf{Grp}}}
\newcommand{\Fld}{{\textbf{Fld}}}
\newcommand{\Sh}{{\textbf{Sh}}}
\newcommand{\Set}{{\textbf{Set}}}
\newcommand{\Mod}{{\textbf{Mod}}}
\newcommand{\PSh}{{\textbf{PSh}}}
\newcommand{\Sch}{{\textbf{Sch}}}
\newcommand{\Top}{{\textbf{Top}}}
\newcommand{\RgSp}{{\textbf{RgSp}}}
\newcommand{\Ring}{{\textbf{Ring}}}


%operatorname -- 3+ letters (lowercase and upper case)
\newcommand{\rad}{{\operatorname{rad}}}
\newcommand{\sep}{{\operatorname{sep}}}
\newcommand{\res}{{\operatorname{res}}}
\newcommand{\sgn}{{\operatorname{sgn}}}
\newcommand{\pre}{{\operatorname{pre}}}
\newcommand{\ord}{{\operatorname{ord}}}
\newcommand{\vol}{{\operatorname{vol}}}
\newcommand{\lcm}{{\operatorname{lcm}}}

\newcommand{\conj}{{\operatorname{conj}}}
\newcommand{\disc}{{\operatorname{disc}}}
\newcommand{\stab}{{\operatorname{stab}}}
\newcommand{\kdim}{{\operatorname{kdim}}}
\newcommand{\diag}{{\operatorname{diag}}}

\newcommand{\trdeg}{\operatorname{trdeg}}
\newcommand{\coker}{{\operatorname{coker}}}
\newcommand{\codim}{{\operatorname{codim}}}

\newcommand{\Ann}{{\operatorname{Ann}}}
\newcommand{\Aut}{{\operatorname{Aut}}}
\newcommand{\Der}{{\operatorname{Der}}}
\newcommand{\Div}{{\operatorname{Div}}}
\newcommand{\Emb}{{\operatorname{Emb}}}
\newcommand{\End}{{\operatorname{End}}}
\newcommand{\Ext}{{\operatorname{Ext}}}
\newcommand{\Gal}{{\operatorname{Gal}}}
\newcommand{\Hom}{{\operatorname{Hom}}}
\newcommand{\Ind}{{\operatorname{Ind}}}
\newcommand{\Inn}{{\operatorname{Inn}}}
\newcommand{\Int}{{\operatorname{int}}}
\newcommand{\Irr}{{\operatorname{Irr}}}
\newcommand{\Jac}{{\operatorname{Jac}}}
\newcommand{\Mor}[1]{\operatorname{Mor}(\textbf{#1})}
\newcommand{\Nil}{{\operatorname{Nil}}}
\newcommand{\Obj}{{\operatorname{Obj}}}
\newcommand{\Out}{{\operatorname{Out}}}
\newcommand{\Pic}{{\operatorname{Pic}\,}}
\newcommand{\Rep}{{\operatorname{Rep}}}
\newcommand{\Res}{{\operatorname{Res}}}
\newcommand{\Spl}{{\operatorname{Spl}}}
\newcommand{\Syl}{{\operatorname{Syl}}}
\newcommand{\Sym}{{\operatorname{Sym}}}
\newcommand{\Tor}{{\operatorname{Tor}}}

\newcommand{\Char}{{\operatorname{char}}}
\newcommand{\Frac}{{\operatorname{Frac}}}
\newcommand{\Frob}{{\operatorname{Frob}}}
\newcommand{\Proj}{{\operatorname{Proj}\,}}
\newcommand{\RMod}{{}_R\operatorname{Mod}}
\newcommand{\SExt}{{\mathcal{E}\emph{xt}}}
\newcommand{\SHom}{{\emph{Hom}}}
\newcommand{\Span}{{\operatorname{span}}}
\newcommand{\Spec}{{\operatorname{Spec}\,}}
\newcommand{\Supp}{{\operatorname{Supp}}}
\newcommand{\Tens}[1]{#1\otimes --}
\newcommand{\fHom}[1]{\operatorname{Hom}(#1, --)}

\newcommand{\mSpec}{{\operatorname{mSpec}}}


% connectors
% \renewcommand{\mid}{\big|}
\newcommand{\sse}{\subseteq}
\newcommand{\subg}{\leqslant}
\newcommand{\supg}{\geqslant}
\newcommand{\ssne}{\subsetneq}
\newcommand{\isosubg}{\lesssim}
\newcommand{\nsse}{\not\subseteq}
\newcommand{\nsubg}{\trianglelefteq}
\newcommand{\nsupg}{\trianglerighteq}
\newcommand{\csubg}{\blacktriangleleft}
\renewcommand{\mid}{\big|}


% Arrows
\newcommand{\rw}{\rightarrow}
\newcommand{\Rw}{\Rightarrow}
\newcommand{\lw}{\leftarrow}
\newcommand{\Lw}{\Leftarrow}
\newcommand{\lrw}{\leftrightarrow}
\newcommand{\LRw}{\Leftrightarrow}
\newcommand{\acton}{\curvearrowright}
\newcommand{\actson}{\curvearrowright}
\newcommand\mapsfrom{\mathrel{\reflectbox{\ensuremath{\mapsto}}}}

% arrows for commutative diagram: create four new angled double-struck arrows
\newcommand\myrot[1]{\mathrel{\rotatebox[origin=c]{#1}{$\Rightarrow$}}}
\newcommand\NEarrow{\myrot{45}}
\newcommand\NWarrow{\myrot{135}}
\newcommand\SWarrow{\myrot{-135}}
\newcommand\SEarrow{\myrot{-45}}


%shorthands
\newcommand{\oln}[1]{{\overline{#1}}}
\renewcommand{\bf}[1]{\textbf{#1}}
\newcommand{\qn}{\quad\newline}


% limits
\newcommand{\Lim}{{\varprojlim}}
\newcommand{\coLim}{{\varinjlim}}
\newcommand{\Dlim}{\lim\limits_{\longrightarrow}}
\newcommand{\coDlim}{\lim\limits_{\longleftarrow}}
\newcommand{\cofHom}[1]{\operatorname{Hom}(--, #1)}


%for saturated ideals in the localization section (I don't like these, I think I'll delete)
\newcommand{\LIm}{{}^e }
\newcommand{\LPre}{{}^c }

%  ╔══════════════════════════════════════════════════════════╗
%  ║                    Custom Commands                       ║
%  ╚══════════════════════════════════════════════════════════╝

% swap varphi and phi
\let\temp\phi
\let\phi\varphi
\let\varphi\temp

\newcommand{\placeholder}{\_\_\_\_\_}

\newcommand{\set}[2]{\left\{#1 \ \middle|\ #2\right\}}

%mod should not have the crazy amount of space to its right!
\renewcommand{\mod}{\,\operatorname{mod}\,}

% dcups
\newcommand{\dcup}{\sqcup}
\newcommand{\Dcup}{\coprod}
% \newcommand{\Dcup}{\bigsqcup}

\newcommand{\nbot}{\mathrel{\rlap{$\bot$}\mkern2mu/}}

%a nice large circle
\newcommand{\tikzcircle}[2][red,fill=black]{\tikz[baseline=-0.5ex]\draw[#1,radius=#2] (0,0) circle ;}%

\makeatletter
\newcommand*\bigcdot{\mathpalette\bigcdot@{.5}}
\newcommand*\bigcdot@[2]{\mathbin{\vcenter{\hbox{\scalebox{#2}{$\m@th#1\bullet$}}}}}
\makeatother

%somehow not a default command
\def\rddots{\mathstrut^{.^{.^{.}}}}

%% new delimiter
\DeclarePairedDelimiter{\ceil}{\lceil}{\rceil}


%  ╔══════════════════════════════════════════════════════════╗
%  ║                     environments                         ║
%  ╚══════════════════════════════════════════════════════════╝


%remember the option: breakable

% create tcolor boxes
\newtcbtheorem[number within=section]{thm}{Theorem}%
{
colback=green!5,
colframe=green!35!black,
fonttitle=\bfseries,
label type=theorem
}{th}

\newtcbtheorem[number within=section, use counter from=thm]{lem}{Lemma}%
{
colback=blue!5,
colframe=black!35!black,
fonttitle=\bfseries,
label type=lemma
}{lm}
\newtcbtheorem[number within=section, use counter from=thm]{prop}{Proposition}%
{
colback=white!5,
colframe=white!35!black,
fonttitle=\bfseries,
label type=proposition
}{pr}
\newtcbtheorem[number within=section, use counter from=thm]{cor}{Corollary}%
{colback=blue!5,
colframe=blue!35!black,
fonttitle=\bfseries,
label type=corollary
}{co}
\newtcbtheorem[number within=section, use counter from=thm]{axiom}{Axiom}%
{colback=black!5,
colframe=yellow!35!black,
fonttitle=\bfseries,
label type=axiom
}{ax}
\newtcbtheorem[number within=section, use counter from=thm]{defn}{Definition}%
{colback=red!5,
colframe=red!35!black,
fonttitle=\bfseries,
label type=definition
}{df}


% for <s-cr> to create \item[$\LRw$] instead of \item
\newenvironment{equivEnumerate}
 {\begin{enumerate}
	}{
\end{enumerate}}

%insure the exercises are properly numbered
\counterwithin{Exercise}{section}
\counterwithin{Answer}{section}


%Custom enviroments
 \newenvironment{note}
 {\begin{tcolorbox}[enhanced, sharp corners, colback=white , borderline={0.2pt}{0pt}{black}]
   \begin{quote}\textbf{Note}:
   }{
   \end{quote}\end{tcolorbox}}

 \newenvironment{intuition}
 {\begin{quote}\textbf{Intuition}:
   }{
   \end{quote}}

\newtcbtheorem[number within=section, use counter from=thm]{titledBox}{}{%
  enhanced,
  sharp corners,
  colback=white,
  colframe=black,
  borderline={0.2pt}{0pt}{black},
  left=2mm,
  right=2mm,
  top=2mm,
  bottom=2mm,
  title style={opacity=1},
  attach title to upper,
  before upper={\textbf{\tcbtitletext}\quad},
  coltitle=black,
  fonttitle=\bfseries,
  separator sign={\quad}
}{box}


%better example and proof environments
\def\exampletext{Example} % If English
\colorlet{colexam}{red!55!black} % Global example color


\newtcbtheorem[number within=section, use counter from=thm]{example}{Example}{% \newtcolorbox[use counter=testexample]{testexamplebox}{%
% Example Frame Start
empty,% Empty previously set parameters
title={\exampletext \thetcbcounter #1},% use \thetcbcounter to access the testexample counter text
% Attaching a box requires an overlay
attach boxed title to top left,
   % Ensures proper line breaking in longer titles
   minipage boxed title,
% (boxed title style requires an overlay)
boxed title style={empty,size=minimal,toprule=0pt,top=4pt,left=3mm,overlay={}},
coltitle=colexam,fonttitle=\bfseries,
before=\par\medskip\noindent,parbox=false,boxsep=0pt,left=3mm,right=0mm,top=2pt,breakable,pad at break=0mm,
   before upper=\csname @totalleftmargin\endcsname0pt, % Use instead of parbox=true. This ensures parskip is inherited by box.
% Handles box when it exists on one page only
overlay unbroken={\draw[colexam,line width=.5pt] ([xshift=-0pt]title.north west) -- ([xshift=-0pt]frame.south west); },
% Handles multipage box: first page
overlay first={\draw[colexam,line width=.5pt] ([xshift=-0pt]title.north west) -- ([xshift=-0pt]frame.south west); },
% Handles multipage box: middle page
overlay middle={\draw[colexam,line width=.5pt] ([xshift=-0pt]frame.north west) -- ([xshift=-0pt]frame.south west); },
% Handles multipage box: last page
overlay last={\draw[colexam,line width=.5pt] ([xshift=-0pt]frame.north west) -- ([xshift=-0pt]frame.south west); }%
}{ex}



\def\prooftext{Proof} % If English
\NewDocumentEnvironment{Proof}{ O{} }
{
    \colorlet{colexam}{black!55!black} % Global example color
    \newtcolorbox[number within=section]{testproofbox}{%
        empty,% Empty previously set parameters
        title={\emph{\prooftext} \thetcbcounter: #1},
        attach boxed title to top left,
        minipage boxed title,
        boxed title style={empty,size=minimal,toprule=0pt,top=4pt,left=3mm,overlay={}},
        coltitle=colexam,
        fonttitle=\bfseries,
        before=\par\medskip\noindent,
        parbox=false,
        boxsep=0pt,
        left=3mm,
        right=0mm,
        top=2pt,
        breakable,
        pad at break=0mm,
        before upper=\csname @totalleftmargin\endcsname0pt,
        % Removed height constraints
        % Added flexible width
        width=\dimexpr\textwidth-3mm\relax,
        % Modified overlays
        overlay unbroken={\draw[colexam,line width=.5pt] ([xshift=-0pt]title.north west) -- ([xshift=-0pt]frame.south west); },
        overlay first={\draw[colexam,line width=.5pt] ([xshift=-0pt]title.north west) -- ([xshift=-0pt]frame.south west); },
        overlay middle={\draw[colexam,line width=.5pt] ([xshift=-0pt]frame.north west) -- ([xshift=-0pt]frame.south west); },
        overlay last={\draw[colexam,line width=.5pt] ([xshift=-0pt]frame.north west) -- ([xshift=-0pt]frame.south west); },%
    }
    \begin{testproofbox}
}
{\end{testproofbox}\endlist}


%  ╔══════════════════════════════════════════════════════════╗
%  ║                    for Dynkin diagrams                   ║
%  ╚══════════════════════════════════════════════════════════╝

\def\row#1/#2!{#1_{\IfStrEq{#2}{}{n}{#2}} & \dynkin{#1}{#2}\\}
\newcommand{\tble}[1]{
   \renewcommand*\do[1]{\row##1!}
   \[
      \begin{array}{ll}\docsvlist{#1}\end{array}
   \]
}

%  ╔══════════════════════════════════════════════════════════╗
%  ║                         links                            ║
%  ╚══════════════════════════════════════════════════════════╝

\definecolor{LinkColour}{HTML}{A64A3B} % favourite
% \definecolor{LinkColour}{HTML}{8C3D32} % 2nd place
\hypersetup{
  colorlinks,
  citecolor=black,
  filecolor=magenta,
  linkcolor=LinkColour,
  urlcolor=blue,
}



%  ╔══════════════════════════════════════════════════════════╗
%  ║                          misc                            ║
%  ╚══════════════════════════════════════════════════════════╝

\stackMath %to be able to stack text in math mode
\pgfplotsset{compat=1.18}
\makeindex

%import options: all, bibliography, book-formatting, booksSetting, customEnvironments, dynkin, hw-formatting, language-setup, newcommands, packages, preamble, proofs, settings, tcolorbox, test.svg,
\begin{document}
\maketitle
\tableofcontents

\newpage
These notes are an extract from EYNTKA Algebra (\cite{NateAlgebra}) that represent a detailed ``preliminary
chapter''. If na\"ive set theory is new to you, this section is going to go through a lot of terminology. The
reader may want to skip to certain sections depending on their background:
\begin{enumerate}
	\item Sets: sets, power-sets, unions and intersections. Logical quantifiers, basic propositions, vacuous
		truth. Induction. Relations, equivalence Relations, partial and total orders.
	\item Functions: injective/surjective/bijective functions and their relation to equivalence
		relations. Image and restriction, composition of functions, associativity of composition, Cardinality
	\item Number theory: GCD and LCM, divisibility, prime factorization, coprime numbers, the euclidean
		algorithm, B\'ezout's Identity, Modular Arithmetics, Fermat's Little Theorem
	\item Graphs: basic properties
	\item Category Theory: intuition, universal property, functors.
\end{enumerate}
You can look over this section if you want a refresher, or skip to section~\ref{sec:equivRelations} to start
from equivalence relations, in particular if the reader is unaware of the relation between equivalence
relations and functions, then this section will be important.

\section{Sets}

As in any field of study, we need to establish a common language for everyone in the field to communicate their ideas.
The accepted language of mathematics today was invented by Georg Cantor in the 19th century and is called \emph{set
theory}. Before Cantor, there's wasn't one universal set of symbols and concepts mathematicians agreed upon to talk
about basic concepts, with different groups of mathematician inventing their own symbology to talk about their ideas.
Nowadays, every field of mathematics grounds their work in set theory \footnote{There is a movement to replace
foundation of mathematics with \emph{type theory} -- for those interested see ref:HERE
} -- algebra is no exception.

We'll be learning a small part of set theory which concerns itself with using the tools of set theory for
talking about concepts. These tools are under the label of \emph{na\"ive set theory}. We will be using these
tools and question little why we are using these particular tools. The reason for the term ``na\"ive'' set
theory is due to history: there were a couple of important paradoxes which showed some inconsistencies in set
theory \footnote{Famously, Russel's paradox showed a logical problem in how we can define sets.  We'll
	actually cover Russel's Paradox in ref:HERE (Perhaps add it to the end of set theory section?) if you're
curious to check it out after reading this section}. From the existence of these paradoxes, a generation of
mathematician’s tried to free set theory from these problems \footnote{one major result of this research was
	the axiomatization  of set theory (i.e. the development of the fundamental rules) by Zermelo and Frankel,
	their set of axioms (i.e. fundamental rules) being called ZF in honour of them. There are a couple of
	axioms that can be added or removed from ZF, but ZF generally forms the core of modern mathematics. If you
keep doing mathematics, you will hear statements like ``This statement is true in ZF''}. Their research
became known as set theory, and the previous understanding of set theory came to be known as \emph{na\"ive
set theory}. Na\"ive set theory wasn't dismissed, instead mathematicians who were interested in making a
solid foundation for mathematics tried figuring out exactly how we should be using na\"ive set theory
without contradiction, while the rest of mathematicians used na\"ive set theory. When we use na\"ive set
theory, we're essentially using the results of [the new] set theory without worrying about the details.
It's like taking advantage of a building without knowing the architectural details of why a support was
put here, or a door was needed there. We will be taking full advantage of this pre-built building for us.
Some footnotes will be added for those who want to peek at the questions that lead mathematicians of the
early 20th century to investigate the details of set theory. For those with a deeper interest in the
foundations of mathematics, a class on logic is recommended, or you can read one of these books:
\footnote{here, some good one's in my Book folder}

%(A jump here when I go into the history) (maybe compare to na\"ive calculus? Using the results without thinking about
%the theory. like understanding the area of a circle is $\pi r^2$, but you don't know why it is $\pi r^2$)


As a consequence of using na\"ive set theory, we'll be defining certain concepts by relying on some [very] basic intuition. For
example, we will assume the idea that we can conceptualize a collection of objects (for example, imagine a bunch
of trees, or a bunch of numbers), or know the basics of numbers (ex. you know how to think of the symbols $1,2,3,...$ as
number, and you can add/subtract/multiply/divide). Definition's after the preliminary chapter will build on the definition's presented
here, and so it's not a prerequisite for chapters after the preliminary chapter to rely heavily on intuition.


\subsection*{Basic Definitions: Sets}



In mathematics, we work with collection of things, and with ``mathematical sentences'' about a collection of things,
where a sentence that must either be true or false. We will start by learning how to describe the collection of things,
and then move on to learning about mathematical sentences. We will make a distinction between a collection of things and
the things inside the collection. We call a collection of things a \emph{set}, and the things inside a set are
\emph{elements}. Another definition of a set is provided by the founder of set theory Georg Cantor

\begin{quoting}
	A set is a Many that allows itself to be thought of as a One
\end{quoting}

In other words, we can think of any arbitrary collection of objects put together as a new single entity. \footnote{Going
deeper than this, asking questions like what are the possible things a set can contain, what are all possible sets, can
there be sets that contain itself, all venture into set theory proper}. The one restriction we give to sets is that all the
elements in a set must be different, in other words, we must have each element to be distinct.

When we write a set, we put curly braces (``$\{$'' and ``$\}$'') around a collection of things (or really, symbols
which could represent things \footnote{Some might be dubious on whether the symbols can represent ``anything''.
For example, is it okay to say that $x$ represents a statement that is both true and false? Is that ``a thing''. This
type of question of what is a ``well-defined'' thing is subject to deeper explorations of set theory}) which are separated by a comma.
\begin{example}{Example of sets}{setsEx}
	\begin{enumerate}
		\item To write a set with the numbers $1,2$ and $3$, we write $\left\{ 1,2,3 \right\}$. The numbers $1$, $2$,
			and $3$ are the elements of $\{1,2,3\}$. Note that $\{1,1,2,3\}$ would not be considered a set
			since the $1$'s repeat. The convention when a set has repeated elements is to have them
			``squished'' into one element to create a set with only distinct elements. This would mean that we would
			consider the set $\{1,1\}$ to be $\{1\}$ where we squished the two $1$'s into a single element. This
			convention also holds if the set has other elements too:
			\[
				\{1,1,2,3\} = \{1,2,3\}
			\]
			If we add an element to a set that's already there, we will also do this ``squishing'' \footnote{Formally,
			(explain what's actually happening here in terms of ZFC?)}

	\item As mentioned, sets contain can anything. The set
		\[
			\left\{ \text{ice-cream}, \text{cheesecake}, \text{chocolate} \right\}
		\]
		is a set that represents $3$ desserts. The set
		%\[
			%\{ \text{Monday, Mardi, \'Sroda, {\dn Bihivaara}, 金曜日}\}
		%\]
		represents the weekdays of the week in 5 different languages.
	\item The set $\{1\}$ is a set with the element $1$. We would not say that $1$ and $\{1\}$ are the same thing, since
		$\{1\}$ is ``a set containing the number one as its element'' and $1$ as ``the number one''. Using mathematical
		notation, $1 \ne \{1\}$
	\item $\{\}$ is also a set, called the empty set. It pops up often enough that we have a special symbol for this
		set: $\emptyset$
	\item Sets can contain sets, so $\{1, \{\smiley, \mathfrak{T}\}, \text{bat}\}$ is a set that contains the element 1,
		the element $ \left\{ \smiley, \mathfrak{T} \right\}$ (which is also a set), and the element bat.
	\item Sets need not have an order, so for example
		\[
			\left\{ \text{ice-cream}, \text{cheesecake}, \text{chocolate} \right\} = \left\{ \text{cheesecake},
			\text{ice-cream}, \text{chocolate} \right\}
		\]
		This also means that a set with days of the week (written in any language) is the same if we re-arrange the days:
		%\[
			%\{ \text{Monday, Mardi, \'Sroda, {\dn Bihivaara}, 金曜日}\} = \{ \text{Mardi, Monday, {\dn Bihivaara}, \'Sroda, , 金曜日}\}
		%\]
\end{enumerate}
\end{example}


Sometimes, we care about the order of elements. If that is the case, we will surround them in parenthesis, ``$($ and
$)$'', and call then $n$-tuples, where $n$ represents the number of elements. For example: $(1,2,3)$ is a $3$-tuple, and
$(1,2,3) \ne (3,2,1)$ even though $\{1,2,3 \} = \{3,2,1\}$ \footnote{In set theory, every mathematical object is
represented as a set. This is no exception in $n$-tuples. A $2$-tuple $(a,b)$ would be $\{a, \{a,b\}\}$. A similar
pattern follows for $n$-tuples}.  $2$-tuples are sometimes called a \emph{pair}. If it is clear from context, it is
a common convention to drop the number in front of the word tuple, i.e., if the writer consistently uses $2$-tuples they
might choose to just call them tuples.


\subsection*{Elements}

As mentioned earlier, If something is inside a set,
we call that an \emph{element} or \emph{object} of the set. We usually denote elements with lower-case letters (ex.
$a,b$)\footnote{Here's another ``beyond the scope of naive set theory'' question: What are the possible elements that
may exist? Question of this varieties lead to logicians like Butrard Russel to come up with \emph{Russel’s Paradox}. This
paradox will actually be covered in the book near the end (see ref:HERE). It's written in such a way that you can skip
to take a look at it if you're curious}. To write ``$a$ in $A$'' in mathematical notation, we use the symbol $\in$.
For example:
\[
	a \in A
\]
is read as ``$a$ in $A$'' (or ``$a$ is in $A$'' to sound more grammatically correct). If we want to say that ``a not in $A$'', we add a slash across our symbol ($/$):
\[
	a \notin A
\]
For example, $1 \in \{1,2,3\}$, but $4\notin \{1,2,3\}$. If we have the empty set $\emptyset$ (i.e. $\{\}$), then it is
always the case that nothing is inside it, i.e., $a\notin \emptyset$ for any $a$. Sets can also contain other sets, so
$\{1\} \in \{\{1\}, 2, \{\text{cheesecake},6\}\}$. Note that $1 \notin \{\{1\}, 2, \{\text{cheesecake},6\}\}$ since the set $\{\{1\}, 2, \{\text{cheesecake},6\}\}$
doesn't contain ``the element $1$'' but the ``the set that contains the element 1''.

\begin{titledBox}{Concept Check}{checkInclusion}
	Is
	\[
		\{1,2,3\} \in \{0, \{\{1,2,3\}\}, 4\}
	\]
	% answer:\footnote{flip this text? No, the set containing the single element$\{1,2,3\}$ is inside the above set}
\end{titledBox}


\subsection*{Equality}

The concept of equality was already used earlier based on our intuition. We will expand a bit further on how to think of
equality in set theory, and we will distinguish two ways of looking at equality:
\begin{gather*}
	A = \{1,2,3\}\\
	A := \{1,2,3\}
\end{gather*}
Starting with the first line, the symbol $=$ means \emph{logically equal}.
Being ``logically equal'' (which we will always abbreviate to saying equal\footnote{This is another big set-theoretic question: what does it really
mean for something to be equal? Bertrand Russel and Alfred Whitehead in their book \emph{Principia Mathematica} famously
took 379 pages before they defined equality. In modern set theory, two sets are equal if they contain the same
elements}) means the two symbols represent the same thing. For example, the number
one-half and two-quarters are equal:
\[
	\frac{1}{2} = \frac{2}{4}
\]
even though the symbol $\frac{1}{2}$ and $\frac{2}{4}$ are different symbols. We would read $\frac{1}{2} = \frac{2}{4}$
as ``one-half is equal to two-quarters''. If we want to say that the two are not equal (i.e. the symbols represent
different things \footnote{This is another question beyond na\"ive set theory: Since all we're using are symbols that
represent objects, where do the objects come from? One answer to this is some universal set with every possible thing as
an element of it, but thinking about what exactly does it mean for a set to contain everything leads to interesting
nuances (one of the problems of a set that contains everything is once again Russel's Paradox)}), we do the same convention as
with $\in$ and put a slash across the symbol: $ \frac{1}{2} \ne  \frac{1}{3}$.

\begin{example}{Examples Of Equality}{examplesEquality}
	\begin{enumerate}
		\item if we want to say two sets have the same elements, that they are each representing the same collection of
			objects, we write $A = B$
		\item $\{1,2,3\} \ne \{2,3,4\}$. Even though both sets have the elements $2$ and $3$, they each contain a number
			the other doesn't have, and so they are not equal
		\item Similarly, $\{1,2\} \ne \{1,2,3\}$ since $3\notin \{1,2\}$
		\item As seen in a elementary class in mathematics, $ \frac{1}{2} = \frac{2}{4}$, but $\frac{1}{2} \ne
			\frac{1}{4}$\footnote{Notice that in all the previous examples, we were comparing sets, while in this
			example we're comparing numbers. In proper set theory, every object is actually a set. Thus, even comparing
			$\frac{1}{2}$ and $\frac{2}{4}$ will be a comparison with sets. We will actually construct $\Q$ and show how
			it's elements are sets in section~\ref{sec:ringOfFracions}}
	\end{enumerate}
\end{example}

The symbol $:=$ means the same thing as $=$, that is, $:=$ is also logically equal. The difference between $:=$ and $=$
is that $:=$ is used to emphasize that the symbol on the left of $:=$ will now be used to represent the thing on the
right of $:=$. This is most often used in definitions to make sure the reader understands that symbol will consistently
be used to represent the same thing many times. Because of this, the symbol $:=$ is read as ``define'' -- the statement
$A := \{1,2,3\}$ can be read that ``Define $A$ to be the set $\{1,2,3\}$''.



\subsection*{Set-Builder Notation}


Many times, we want to have a set of elements that satisfy only certain properties. For example, from the set
of all deserts, I only want the ones that contain ice-cream, or from the set of all numbers I want only the
prime numbers, or from the set of all months, I want the months that end in ``ber''. The way we'll capture
this idea with sets is by using \emph{set builder-notation} to build a new set given an old one. If $S$ is a
set and $P$ is some sentence describing a property (ex. being an even number, or being a prime number, or
representing ``this desert has ice-cream''), then we'd write


\begin{gather*}
	A = \{s \in S\ |\ P\}\\
	\text{or}\\
	A = \{s \in S\ :\ P\}\\
\end{gather*}
and is read ``s in S such that [s satisfies the property P]''. In the next section, we shall be more precise
about what type of sentence can be used:

\begin{example}{Set Builder Notation}{setbuilderNotation}
	\begin{enumerate}
		\item If
			\begin{gather*}
				M = \{ \text{Janurary},\ \text{February},\ \text{March},\ \text{April},\ \text{May},\ \text{June},\\
					\text{July},\ \text{August},\ \text{September}, \text{October},\ \text{November},\ \text{December}\}
			\end{gather*}
			and $D = \set{m \in M}{\text{m ends in ``ber''}}$, then
			\[
				D = \{\text{September, October, November, December}\}
			\]
		\item If $W = \{\text{hello, (words in other languages)}\} $, and $V = \set{w \in W}{\text{$w$ is in
			Vietnamese}}$, then
			\[
				V = \{\text{here}\}
			\]
		\item if $\N = \{0,1,2,3,\cdots\}$ and $P = \set{n \in \N}{\text{$n$ is a prime number}}$ then
			\[
				P = \{2,3,5,7,11,\cdots\}
			\]
			Note that $P' = \set{n \in \N}{\text{the only numbers that divide $n$ are $1$ and $n$}}$ is the exact same
			set (i.e. $P = P'$). Different properties can produce the same sets.
		\item if $\N = \{0,1,2,3,\cdots\}$ represents the
set of all non-negative numbers with no decimal places, the set
\[
	A = \{x \in \N\ | x = 2y,\ y \in \N\}
\]
would represent all even numbers \footnote{This notation has it's limitations if you dig deeper in it. You can
create some very interesting paradoxes (Russel's paradox being one of them). We will not construct such contradictory
cases in the book}, and can be read ``$x$ in the natural numbers such that $x$ can be twice any natural number''
	\end{enumerate}

\end{example}

\subsection{Propositions and Predicates}\label{sec:propositions}

We will now explore how to formulate questions and mathematical statements using mathematical language.

Intuitively, there is some understanding of the meaning of a sentence\footnote{The nature of how we derive
meaning from sentences is an interesting logical and philosophical debate that goes beyond the scope of
this book. See Russel, Wittgenstein, \placeholder, for more details} . If I say ``I think blue is pretty''
then in your mind some information got passed on. If we make a sentence that has a truth value associated
to it (the sentence can either be true or false) \footnote{Development in a field called \emph{topoi
theory} are trying to generalize our notion of proposition beyond true and false, an interesting
developing field to look into} , like ``The number $2$ is prime'', then  we call such a sentence a
\emph{proposition}. If we have a sentence that takes in a variable that will make it depend on its truth
value, for example ``The number $x$ is a prime number'', then this will be called a \emph{predicate}.
Propositions and predicates are at the heart of mathematical communication. We have already encountered
predicates before when we used set-builder notations: the restriction \emph{has to be a predicate}. For
convenience, we assign a symbol to a proposition. For example, we can have $P = \text{``2 is a prime
number''}$. Then $P$ would be true. However, $Q = \text{``2 is not a prime number''}$ is a false. For
predicates, we shall usually assign a symbol and a variables, for example:
\[
	R(x) = \text{``$x$ is a prime number'''}
\]
where $x$ may be any natural number, then this predicate truth-value (that is, whether it is true or false) depends on the value the variables $x$ will take.

% (don't know where to move these yet)
% (Maybe add somewhere how $a \in A$ can be thought of as the sentence $\exists x \in A \text{ such that } x = a$ is true?)
% (explain why it tries to get rid of ambiguity. For example, if we do not carefully define proposition, have the
% proposition $P$ such that $P$ is an interesting number. What does ``interesting number'' mean? )


\begin{example}{Predicates}{propositionsEx}
	\begin{enumerate}
		\item let $P(x) =  $``$x$ is a month of the year''. Then
			\[
				P(\text{September}) = \text{true} \qquad P(\text{Wednesday}) = \text{false} \qquad P(\text{Laughing})
				= \text{false}
			\]
		\item The variable $x$ can also be other propositions. For example we can have $P(x) = $``$x$ is a dog'' and
			$Q(x) =$``$x$ is an animal''.
		\item  Propositions can be composed, for example consider $R(x) =$``$P(x)$, then $Q(x)$'' where $P(x),
			Q(x)$ are the propositions from the above. When is this proposition true?
		\item Proposition can also take in multiple variables, for example
			\[
				P(x,y) = (x < y)
			\]
	\end{enumerate}
\end{example}

Predicates of the form
\begin{center}
	if $P(x)$ then $Q(x)$
\end{center}
are some of the most common, allowing for relations being made between
different concepts. The symbol used to symbolize an ``if, then'' relation is $\Rw$, so
\[
	\text{If $P(x)$ then $(Q)$} \qquad \text{becomes}\qquad P(x) \Rw Q(x)
\]
Which is usually read ``$P(x)$ implies $Q(x)$''. A  common mistake is to assume that if $P(x) \Rw Q(x)$, then $Q(x) \Rw
P(x)$. This need not be the case. For example, define the following two predicates
\[
	A(x) = (x\mid 2) \Rw (x\mid 4) \qquad A'(x) = (x\mid 4) \Rw (x \mid 2)
\]
When is predicate $A(x)$ true? When is $A'(x)$ true? If it is true that $P(x) \Rw Q(x)$ always implies $Q(x)
\Rw P(x)$ we use the symbol $P(x) \LRw Q(x)$ and read $\LRw$ as ``if and only if'' (sometimes abbreviated to
iff).


\subsection*{Vacuous Truth}

Let's say we had two proposition $P(x,y)  = x \Rw y$, and let's say that $x$ was false, while $y$ is true. What is the
truth value of $P(x,y)$? For example, let's say that
\begin{align*}
	x &= \text{I'm walking in the rain}\\
	y &= \text{I am wet}
\end{align*}

Then if you are \emph{not} walking in the rain, but \emph{are} wet, then what should be the value of $P(x,y)$?
Perhaps more interesting is the following example:
\begin{enumerate}
	\item Let's say there are no cell-phones in a room. Then is the statement ``all phone are turned off'' true?

	\item A famous such example is the following, is the statement ``The king of France is not wearing pants'' true
		or false if France does not have a king?
\end{enumerate}

There are many solutions to this problem. The mathematically accepted solution to this problem is something
called \emph{vacuous truth}, which gives the following resolution:

\begin{center}
	$(\text{False} \Rw \text{True})$ is True, and $(\text{False} \Rw \text{False})$ is \emph{also} True
\end{center}

The idea being that any ``nonsense'' statement is always ``true'' in the most trivial way. This interpretation
of truth is not a unique one, but it is one that is very convenient as shall be seen by the reader.

\subsection*{Logical Quantifiers}

We will now explore how to create more complicated properties by adding to our tool belt $6$ ``words'' from which we can
construct mathematical sentences:
\begin{defn}{Quantifiers And Connectives}{quantifiersAndConnectives}
	We define the following $3$ quantifiers to be used in propositions:
	\begin{enumerate}
		\item $\forall$: the for-all quantifier.
		\item $\exists$: the \emph{existence} quantifier
		\item $\exists!$ the \emph{unique existence quantifier}
	\end{enumerate}
	and the following 4 connectives to be used for propositions
	\begin{enumerate}
		\item $\vee$: disjunction (or) connective. If $P, Q$ are propositions, then $P\vee Q$ mean that if
			one of or both of $P$ or $Q$ is true, then the new poposition $P\vee Q$ is true.
		\item $\wedge$: conjunction (and) connective. If $P, Q$ are propositions, then $P\wedge Q$ mean that if
			 both $P$ and $Q$ are true, then the new proposition $P\wedge Q$ is true. If one of them are
			 false, then $P\wedge Q$ is false.
		 \item $\Rightarrow$: implication connective: $P \Rw Q$ is false if $P$ is true and $Q$ is false. In
			 other words, if $P$ is true, then $P$ implies $Q$ means $Q$ is true.
		 \item $\neg$ negation (not) connective. If $P$ is a proposition, then if $P$ is true then $\neg P$ is
			 false, and similarly if $P$ is false then $\neg P$ is true.
	\end{enumerate}
\end{defn}

The keen reader may ask if there are any other quantifiers or connectives, and indeed there can be many more.
However, it is in fact sufficient to define these connectives/quantifiers to completely describe any
[first-order] mathematical sentence\footnote{This would be explored in a course on logic}. Using these
symbols, can the reader translate some of the propositions earlier in the book using these symbols?

Recall that we have defined the symbol $\in$ earlier to specify an element belongs to a set. This is a special
symbol which is neither a quantifier nor a connective. It is considered a special symbol that is part of the
language of set theory.

The order of quantifiers is important. For example:
\[
	\forall x,\ \exists y,\ x < y
\]
means that for every $x$, there exists a $y$ that is greater than it, which is certainly true when working
with numbers. However:
\[
	\exists x,\ \forall y,\ x < y
\]
means that there exists an $x$ that is smaller than all $y$, which is certainly not possible since we can
always find a smaller number.

With these, we can define certain operations we can perform on sets (don't forget to add visuals):
\index{Set Operations}
\begin{defn}{Set Operations}{setOperations}
		\begin{enumerate}
		\item Define a subset to be:
			\[
				A \sse B \LRw (x \in A \Rw x \in B)
			\]
			Sometimes the symbol $A \subset B$ is also used. Since it is possible that $A \subset B$ $A = B$ (you can check
			this by the definition), the symbol $\sse$ will predominantly be used to emphasize this relation. If we want to make it
			clear that $A \ne B$, we write $A \subsetneq B$, we will write $A \ssne B$.

			\qn
			For any set $S$, we always have $\emptyset \sse S$ and $S \sse S$.  With this notation, we also have that $A
			\sse B$ and $B \sse A$ $\Rw$ $A = B$
		\item Define the intersection of two sets to be:
			\[
				A\cap B := \set{x}{x \in A\, \text{ and } x \in B}
			\]
		\item Define the union of two sets to be:
			\[
				A\cup B := \set{x}{x \in A \text{ or } x \in B}
			\]
			if $A \cap B = \emptyset$, we say that $A$ and $B$ are \emph{disjoint}
		\item Define difference between sets to be:
			\[
				A\setminus B := \set{x}{x \in A\, \text{ and } x\notin B}
			\]
			sometimes also written as $A - B$. If $T \sse S$, then we say that $S \setminus T$ is the \emph{compliment} of
			$T$, usually written $T^c$.
	\end{enumerate}
\end{defn}

Another very important definition

\begin{defn}{Product Of Sets}{prodSets}
	Let $A$ and $B$ be sets. Then
	\[
		A\times B = \set{(a,b)}{a \in A,\ b \in B}
	\]
	is the set of all $2$-tuples (called just tuples for short).  This construction can be extended
	indefinitely:
	\[
		X_1\times X_2\times \cdots = \set{(x_1, x_2, ...)}{x_1 \in X_1,\ x_2 \in X_2,\ ...}
	\]
	This notation is also abbreviated as:
	\[
		\prod_{i=1}^n X_i
	\]
\end{defn}

For example:

\begin{example}{Product Of Sets}{prodSets}
\[
	\{1,2\} \times \{a,b\} = \{(1,a), (2,a), (1,b), (2,b)\}
\]
\begin{center}
	and
\end{center}
\[
	\{a,b\} \times \{1,2\} = \{(a,1), (b,1), (a,2), (b,2)\}
\]
\end{example}

(A word on arbitrary products? i.e. $\prod_{i \in I} X_i$)

\begin{Exercise}[label=ex:basicSetExercises]
	\Question Let $\{E_n\}_{n=1}^\infty$ be a collection of sets indexed by the natural numbers. Define
  \[
  	\limsup\{E_n\}_{n=1}^\infty = \bigcap_{k=1}^\infty \bigcup_{n=k}^\infty E_n \qquad \liminf \{E_n\}_{n=1}^\infty = \bigcup_{k=1}^\infty \bigcap_{n=k}^\infty E_n
  \]
  show that
  \begin{gather*}
  	\limsup \{E_n\}_{n=1}^\infty = \set{x}{x \in E_N\text{ for infinitely many $n$}}\\
  	\liminf \{E_n\}_{n=1}^\infty = \set{x}{x \in E_N\text{ for all but finitely many $n$}}\\
  \end{gather*}

  \Question (Demorgans La)
\end{Exercise}
\begin{Answer}[ref={ex:basicSetExercises}]
  \Question
\end{Answer}


\subsection*{Naming the Numbers and Common Sets}

With the basics of set down, we can introduce some of the most common sets and their notation
\begin{defn}{Common Sets}{commonSets}
	\begin{enumerate}
		\item As we've already seen, if a set has no elements, the symbol $\emptyset$ represents the empty-set
		\item If a set has a single element, we call that set a \emph{singleton}. For example, $\left\{ 1 \right\}$ and
			$\left\{ \pi \right\}$ are singletons.
		\item If $X$ is a set, $P(X)$ is the set of all subsets of $X$\footnote{The existence of this set is an axiom
			in ZFC}. In set builder notation, it can be denoted $P(X) = \set{S}{S \sse X}$. For example, if $X = \{a,b,c\}$, then
			\[
				P(X) = \{ \emptyset, \{a\}, \{b\}, \{c\}, \{a,b\}, \{a,c\}, \{b,c\}, \{a,b,c\}\}
			\]
		\end{enumerate}
\end{defn}

\begin{defn}[breakable]{Sets Representing Numbers}{setsRepNumbers}
	\begin{enumerate}
		\item $\N := \{0,1,2,3,\cdots\}$ is the set of all \emph{natural numbers} (some call it the \emph{counting numbers}
			or \emph{whole numbers} ). Natural numbers are all positive
			numbers which do not have any decimal places. Sometimes, people define $\N$ to not have $0$. We will always have
			$\N$ contain the number $0$ (for reasons which will be clear much later in the book \footnote{See
				section~\ref{sec:monoids}}). If we want only the [strictly] positive number \footnote{There is no complete agreement on whether
			$0$ is positive or negative: $0$ is either both positive and
		negative or neither positive nor negative. It usually comes down to the purposes of the context. If we allow $0$
	to be a positive number, we sometimes say that all numbers larger than $0$ are \emph{strictly positive}}, we will
	write $\N_{>0}$.
		\item $\Z := \{\cdots -3,-2,-1,0,1,2,3,\cdots\}$ is the set of all \emph{integers}. Integers are the natural
			numbers along with their negative counter-part. Sometimes, the natural numbers are called the non-negative
			integers.
		\item $\Q := \set{ \frac{p}{q} }{ p,q \in \Z,\ q\ne 0 }$ is the set of all \emph{rational numbers}. For example,
			$\frac{1}{2} \in \Q$. The symbol
			$\Q$ is used since numbers of the form $\frac{p}{q}$ are called \emph{\textbf{q}uotients}. Notice also that
			this is one of the instances where we use the convention that if there is a repeated element, we ``squish''
			to a single element. If we did not take this convention, this set wouldn't be defined, since there would be
			many repeated numbers (for example $ \frac{1}{2} = \frac{2}{4} = \frac{4}{8}$ and so forth)
		\item $\R$ represents any number with any decimal expansion (ex. $\pi, e, 1, \frac{1}{3}$), i.e.,
			all the numbers on the number line. If you've learned about Complex numbers or
			Quaternions in high school, these numbers are not real numbers, and will be defined and studied later in this
			book. The set $\R$ is called the \emph{real numbers}. We can abuse notation (meaning use the notation in
			a way bends the rules a little) and write $\R$ as
			\[
				\R := \set{\cdots D_3D_2D_1.d_1d_2d_3\cdots}{D_i, d_i \in \Z}
			\]
			where $\cdots D_3D_2D_1.d_1d_2d_3\cdots$ represents some arbitrary string of numbers put together where there
			can only be finitely many $D_i$'s, but there can be infinitely many $d_i$'s. In other word, the number
			\[
				0.123123123123\cdots
			\]
			where $123$ continues on forever on the right is in the set $\R$, but
			\[
				\cdots123123123
			\]
			where $123$ continues on forever on the left is not in the set $\R$.
	\end{enumerate}
\end{defn}

Some find it intuitive to think about $\R$ has having ``no gaps'', which means we cannot perfectly split $\R$ into two
using the less than and greater than operator ($<$ and $>$) without missing a number. For an example of a split, we
can split $\N$  into $\N_{ \frac{5}{2}^+} = \set{x \in \N}{x > \frac{5}{2}}$ and $\N_{ \frac{5}{2}^-} = \set{x \in \N}{x
< \frac{5}{2}}$, and we have successfully split $\N$ into two without missing a number (i.e. every natural number is
either in $\N_{ \frac{5}{2}^+}$ or $\N_{ \frac{5}{2}^-}$). If we chose to split around $5$, then neither the set
$\N_{5^+}$ nor the set $\N_{5^-}$ would contain $5$, meaning it does not split $\N$. With the Real numbers, there does
not exist a perfect split. You might ask if $\Q$ has this same property, but we will prove that we can in fact always
split it: We will show $\Q_{\sqrt{2}^+} = \set{x \in \Q}{x > \sqrt{2}}$ and $\Q_{\sqrt{2}^-} = \set{x \in \Q}{x
< \sqrt{2}}$ splits $\Q$ (see ref:HERE).

The fact that $\R$ cannot be split into two using $>$ and $<$ is actually at the heart of many formal definitions of  $\R$
which is covered in an Analysis I class. In this book, we will not need to think of $\R$ in this fashion, and so we'll stick with just
using intuition for now. \footnote{Google \emph{Dedikend Cut construction} or \emph{Completion of Rational Numbers} if
you're curious to see a formal definition of $\R$}


\begin{center}
	\url{https://math.stackexchange.com/questions/2834876/convergence-of-geometric-series-with-r1}
\end{center}
\subsection{Induction}

Here. The whole $(p(x) \Rw p(x+1)) \Rw \text{here}$ this as the formal definition? + some intuition?

\subsection{*Pigeon-hole Principle}\label{sec:pigeonholePrinciple}

(Comes up in algeraic number theory, so should add it here)

\subsection{Relations and Equivalence Relations}\label{sec:equivRelations}

What does it mean two things to be equal?  There
are actually many ways, depending on what you want to be ``equal''. For example,
\begin{enumerate}
	\item we may say that two people have the same birthday, even if they are not born in the same year.
	\item \textcolor{red}{more examples her}
\end{enumerate}

In this section, we shall explore how to capture this more general notion of equality.

\subsection*{Relations}

We start of with the more general notion of \emph{relation}.  Many times when working with sets, some subsets
can be thought of as having certain properties for example, we can split the set $\Z$ into even and odd
numbers, or we have some numbers being bigger than others, or that some numbers are multiples of another. The
following is this idea is formalized:

\index{Relation}
\begin{defn}{[Binary] Relation}{relation}
	Let $X$ be a set. Then any collection of pairs $(x,y)$, $x,y \in X$ is called a [binary] \emph{relation}. If $(a,b) \in R$, we will
	say that \emph{x is related to y}. If $(a,b) \in R$, we will write $aRb$.

	\qn
	More generally, we can increase the number of times we take the product (i.e. $n > 2$)
	(ex. $X \times X \times \cdots X)$. This would be called an \emph{n-relation}.
\end{defn}

We shall almost always think of relations

\begin{example}{Relations}{relationsEx}
	\begin{enumerate}
		\item Take $\Z$ and define the set $<$ to be
			\[
				(x,y) \in < \text{ if and only if $x$ is strictly less than $y$}
			\]
			Then $(3,4) \in <$, $(3,3) \notin <$, and $(4,3) \notin <$. If $(a,b) \in <$, we would usually
			write:
			\[
				a < b
			\]
			and if $(a,b) \notin <$, we would write
			\[
				a \not< b
			\]
			A similar relation is the $\le$ relation, where we also allow for two numbers to be equal
		\item Take $\Z$ and define the set $R_{\text{div}}$ to be:
			\[
				(x,y) \in R_{\text{div}} \text{ if and only if $x$ divides $y$}
			\]
			Then $(2,4) \in R_{\text{div}}$. What other elements are in this set? This set is usually denoted
			$|$ so that if $a$ divides $b$ then:
			\[
				a | b
			\]
		\item for a non-mathematical example, if we let $X$ be the set of all humans, then we may define
			a relation ``is the ancestor of''. Another similar but different relation is ``is the parent of''
			where two humans are related is one is the parent of another.
	\end{enumerate}
\end{example}

A common relation to put on objects give a sense ordering:
\index{Order on Set}
\begin{defn}{Order}{orderOnSet}
	Let $R$ be a relation on $X$. Then $R$ is a \emph{partial order} if it is:
	\begin{enumerate}
		\item \textbf{reflective}: For all $x \in X$, $xRx$ ($x$ always relates to itself)
		\item \textbf{antisymmetric}: if $xRy$ and $yRx$, then $x=y$\footnote{a relation is symmetric if $xRy$
			then $yRx$. Hence, antisymmetry means no elements can be symmetric with another elements besides
		itself}
		\item \textbf{Transitive}: if $xRy$ and $yRz$, then $xRz$.
	\end{enumerate}
	Furthermore, a partial order is a \emph{total order} (or simply an \emph{order}) if it is furthermore:
	\begin{enumerate}
		\item[4.] \textbf{Connected}: For each $x,y \in X$, either $xRy$ or $yRx$
	\end{enumerate}
\end{defn}

The canonical total order is $\le$ on numbers, and the canonical partial order is $\sse$ on sets.

\begin{example}{Orders}{ordersEx}

	\begin{enumerate}
		\item Let $\N^2 = \N \times \N$. Give some examples on $\N^2$ of different ordering (Ex. dictionary, and others)
		\item Take any dictionary. Then they are lexicographically ordered (more words here)
		\item All numbers we introduced so far come have an order defined: $(\N, \le), (\Z, \le), (\Q, \le), (\R, \le)$.
			This is the order define in high school (so $3 \le 4$ or $3 \le \pi$)
		\item other cool example
	\end{enumerate}
\end{example}

\subsection*{Equivalence Relations}

The most common type of relation that we shall see is the equivalence relation which generalizes the notion of
equality:
\index{Equivalence Relation}
\begin{defn}{Equivalence Relation}{equivalence Relation}
	Let $S$ be a set and $x,y,z \in S$ be elements of $S$. An \emph{equivalence relation} on a set $S$ is relation $\sim$ satisfying
	\begin{enumerate}
		\item \textbf{(reflective)} $x\sim x$ (i.e. every element must be related to itself)
		\item \textbf{(symmetric)} if $x\sim y$ then $y\sim x$
		\item \textbf{(Transitive)} if $x\sim y$ and $y\sim z$ then $x \sim z$
	\end{enumerate}
\end{defn}

\begin{example}{Equivalence Relations}{equivalenceRelationsEx}
	\begin{enumerate}
		\item Equality $(=)$ is the prototypical example of an equivalence relation. For example, If we have $a,b,c \in S$
			with $aRb \LRw a = b$, then
			\begin{enumerate}
				\item $a = a$
				\item $a =b \Rw b = a$
				\item If $a = b$ and $b = c$, then $a = c$.
			\end{enumerate}
			In a ways, an equivalence relation is a generalization of equality.
		\item Let $X$ be any non-empty set. Then we may make all elements equivalent to each other.
			Similarly, we may make all elements only equivalent to themselves.
		\item Let $P$ be the set of all people. Check that ``has the same birthday'' is an equivalence
			relation
		\item Two equivalence relations from high school that the reader may have encountered is with
			triangles. Check that being \emph{similarto} and being \emph{congruent to} are equivalence
			relations.
		\item Define $\sim$ on $\Z$ to be $a \sim b$ if and only if $|a| = |b|$. Then $\sim$ is an equivalence
			relation.
		\item Is $\le$ an equivalence relation?
		\item Let $R$ be the empty relationship on any non-empty set. Is $R$ an equivalence relation?
	\end{enumerate}
\end{example}

Given an equivalence relation $\sim$ on $S$, we may split $S$ into subsets defined as:
\[
	(S/\sim) = \set{T \sse S}{x,y \in T \iff x\sim y}
\]
For example, if $S = \{1,2,3,4,5,6,7,8,9\}$, and:
\[
	1\sim 2 \qquad 3\sim 4\sim 5 \sim 6 \qquad 7\sim8\sim 9
\]
Then:
\[
	S/\sim = \{\{1,2\}, \{3,4,5,6\}, \{7,8,9\}\}
\]
In fact, we may define an equivalence relation on a set $S$ by breaking it up into disjoint sets who union of
elements would be $S$. This is called a partition:

\index{Partition}
\begin{defn}{Partition}{partition}
	Let $S$ be a set. Then a \emph{partition} of $S$ is a collection of subsets of $S$ such that
	\begin{enumerate}
		\item for each subset $S_i, S_j$, $S_i\cap S_J = \emptyset$
		\item $\cup_i S_i = S$
	\end{enumerate}
	each partition is called an \emph{equivalence class}.
\end{defn}


\subsection{Functions}\label{sec:functions}

In mathematics, it is very desirable to relate one set to another set, or said differently how two sets may
``interact''. For example, it may be desirable that
to each positive natural number $\N_{>0}$, we associate to it the number of prime factors it has, for example:
\begin{align*}
	1 &\mapsto 1\\
	2 &\mapsto 1\\
	3 &\mapsto 1\\
	4 &\mapsto 2\\
	5 &\mapsto 1\\
	6 &\mapsto 2\\
	  &\vdots \\
	16&\mapsto 4\\
	  &\vdots
\end{align*}

Another more abstract relation we may ask for is given a set of names, say
\[
	\left\{\text{Alfred},
	\text{Socrates}, \text{Buddha}, \text{Oscar Petesron}\right\}
\]
the function returns a set
containing some information about that person:
\begin{align*}
	\text{Alfred} &\mapsto \{\text{Batman's Butler}, \text{Friend of Bruce}\}\\
	\text{Socrates} &\mapsto \{\text{Philosopher}, \text{lived in Athens}, \text{Never wrote his
	work down}\}\\
			\text{Buddha} &\mapsto \left\{\text{Real name Siddhartha Gautama},
	\text{Starter of a Religion}, \text{Born in Lumbini}\right\}\\
				\text{Oscar Peterson} &\mapsto \left\{\text{Greatest Jazz Pianist},
		\text{Canadian}, \text{``Maharaja of the keyboard''}, \text{Lived to 82} \right\}
\end{align*}

As a final example, we may imagine a every location of a map of the world being represented as an element in a set. Then we may relate a point on the map to
the temperature at that location. Generalizing this pattern, given some set $X$, we may want to assign it an
element from another set $Y$.
\index{Function}
\begin{defn}{[set] Functions}{function}
	Let $X$ and $Y$ be sets. Then a \emph{function} is a subset $\Gamma(f) \sse \set{(x,y)}{x \in X,\ y \in Y}$ such that
	\begin{enumerate}
		\item \textbf{Total} or \textbf{Serial} (left-total): For every $x \in X$, there exists an element $(x,y) \in \Gamma(f)$
		\item \textbf{Functional} (right-unique): Every $x \in X$ is associated with one and only one $y \in Y$, that is, it cannot be that $(x,y), (x,y') \in
			\Gamma$, for $y \ne y'$ where $y,y' \in Y$
	\end{enumerate}
	This is usually written as:
	\[
		f: X \to Y \qquad f(x) = y \qquad \text{or} \qquad X \xrightarrow{f} Y
	\]
	where $f$ is called the \emph{function} and $\Gamma(f)$ is called the graph of the function and is defined
	to be:
	\[
		\Gamma(f) = \set{(x,y)}{x \in X,\ y \in Y,\ f(x) = y}
	\]
	If $f: X \to Y$, $X$ is called the \emph{domain} and $Y$ is called the \emph{codomain}. The collection of
	all functions from two sets $X,Y$ is denoted:
	\[
		\Hom_{\textbf{Set}}(X,Y) \qquad \text{or} \qquad Y^X\ \ ^(\footnote{Notice this is a set operation. This
		particular notation shall be elaborated on with more category theory in section ref:HERE}^)
	\]
\end{defn}

Function are also sometime called maps or mappings. Sometimes, we want to define a function, but don't need to
reference to it again (in computer science, this is called an \emph{anonymous function}). In this case, we may
use the following notation:
\[
	x \mapsto y
\]
instead of $f(x) = y$. If $X = Y$ so that $f: X \to X$, then we may think of $f$ as a relation which has the
property of being functional and total. If more than one function with the same domain and codomain is
defined (say $f, g$ with domain $X$ and codomain $Y$), then the declaration is often shorthanded to $f,g:
X \to Y$.

\begin{example}{Functions}{FunctionEx}
	\begin{enumerate}
		\item The first three examples given are examples of functions, since every element in the domain
			get's mapped to a unique element in the codomain.
		\item Some other perhaps more familiar functions to the reader are
			\begin{enumerate}
				\item $x \mapsto e^x$, i.e. $f: \R \rw \R,\ f(x) = e^x$
				\item $\log: \R_{>0} \rw \R$, $x \mapsto \log(x)$
				\item $p(x): \R \rw \R$, $p(x) = ax^2 + bx + c$, for $a,b,c \in \R$
			\end{enumerate}
		\item Let $X$ be an arbitrary non-empty set. Then there always exists at least one function from $X$
			to $X$, namely:
			\[
				\id_X: X \to X \qquad \id_X(x) = x
			\]
			that is, $\id_X$ maps every element to itself.
	\end{enumerate}
\end{example}

\begin{defn}{Image and Restriction of Function}{imageOfFunc}
	\begin{enumerate}
		\item 	Let $S \sse X$ and $f: X \to Y$. Then:
	\[
		f(S) := \set{f(x)}{x \in S}
	\]
		this is called the \emph{image} of $S$. If $S = X$, then $f(X)$ is called the \emph{image} or \emph{range}
		of the function $f$, and is commonly denoted $\im(f)$.
	\item Let $S \sse X$. Then $f|_S$ is a function $f|_S: S \to Y$ such that
		\[
			f|_S(x) = f(x)
		\]
		that is, $f|_S$ is the \emph{restriciton} of $f$ by $S$.
	\end{enumerate}
\end{defn}

As a function is just some ``syntactic sugar'' to more easily represent the set $\Gamma_f$, two functions
$f,g$ are equal if $\Gamma_f = \Gamma_g$. In terms of function notation this means that $f,g$ are equal if for
all elements $x$ in the domain:
\[
	f(x) = g(x)
\]
An advantage of the left-total and right-unique property is that it is easy to compose functions together,
that is it is easy to chain together the output of one function to be the input of the next:

\begin{prop}{Composition of Functions}{funcComposition}
	Let $f: X \rw Y$, $g: Y \rw Z$, $h: Z \rw W$ be function. Define $(g\circ f)$ to be $(g\circ f)(x) =
	g(f(x))$. Then $(g\circ f)$ is a function and is called \emph{g compose f}. The process is called
	\emph{function composition}.
\end{prop}

\begin{Proof}
	Make sure that $g\circ f$ satisfies the properties of functions
\end{Proof}


\begin{example}{Function Composition}{funcCompEx}
	\begin{enumerate}
		\item Let $f: \R \to \R$, $f(x) = x^2$ and $g: \R \to \R^2$ be $g(x) = (x,0)$. Then:
			\[
				(g\circ f)(x) = (x^2,0)
			\]
			A function may also be composed with itself, given the domain and codomain match:
			\[
				(f\circ f)(x) = x^4
			\]
		\item Let $f: X \to Y$ be any function, $\id_X, \id_Y$ the identity functions defined in
			example~\ref{ex:FunctionEx}. Then:
			\[
				f\circ \id_X = f \qquad \id_Y \circ f = f
			\]
	\end{enumerate}
\end{example}

This chaining together is in a way ``state-independent'': given three functions $f,g,h$, the function $(g\circ
f)$ following by composing with $h$ is equal to if we first just do $f$, then follow by taking the combined
state $(h\circ g)$\footnote{can you think of an example where this is not the case in your everyday life?
This would greatly help in understanding associativity}:

\begin{prop}{Associativity of Functions}{associativityFunctions}
	Let $f: X \rw Y$, $g: Y \rw Z$, $h: Z \rw W$ be function. Then:
	\[
		h\circ(g\circ f) = (h\circ g)\circ f
	\]
	Pictorially, function composition may be represented by the following diagrams:
	\[
		\begin{tikzcd}
A \arrow[r, "f"] \arrow[rr, "g\circ f", bend left] & B \arrow[r, "g"] & C
\end{tikzcd}
\qquad \text{or} \qquad
\begin{tikzcd}
A \arrow[r, "f"] \arrow[rd, "g\circ f"] & B \arrow[d, "g"] \\
                                        & C
\end{tikzcd}
\qquad \text{or} \qquad
\begin{tikzcd}
& C                 \\
A \arrow[r, "f"] \arrow[ru, "g\circ f"] & B \arrow[u, "g"']
\end{tikzcd}
	\]
	and associativity may be drawn as:
	\[
	\begin{tikzcd}
A \arrow[r, "f"] \arrow[rr, "g\circ f", bend right] & B \arrow[r, "g"] \arrow[rr, "h\circ g", bend left] & C \arrow[r, "h"] & D
	\end{tikzcd}
\]
\end{prop}

\begin{Proof}
	This is a long and tedious proof that the reader should attempt at least once in their life.
\end{Proof}


These diagrams have the following special property: start on one of sets in the diagram (say $A$) and travel
down the arrows to reach another set (say to $D$ in the last diagram). Now start again at $A$ and travel down
to $D$ via a different path in the diagram. Then notice that the compositions of functions agree with each
other in the two paths. For example, take the following path:
\[\begin{tikzcd}
	A & B & C & D
	\arrow["f", color={rgb,255:red,214;green,92;blue,92}, from=1-1, to=1-2]
	\arrow["g", from=1-2, to=1-3]
	\arrow["h", from=1-3, to=1-4]
	\arrow["{g\circ f}"', bend right, from=1-1, to=1-3]
	\arrow["{h\circ g}", color={rgb,255:red,214;green,92;blue,92}, bend left, from=1-2, to=1-4]
\end{tikzcd} \qquad \text{and} \qquad
\begin{tikzcd}
	A & B & C & D
	\arrow["f", from=1-1, to=1-2]
	\arrow["g", from=1-2, to=1-3]
	\arrow["h", color={rgb,255:red,214;green,92;blue,92}, from=1-3, to=1-4]
	\arrow["{g\circ f}"', color={rgb,255:red,214;green,92;blue,92}, bend right, from=1-1, to=1-3]
	\arrow["{h\circ g}", bend left, from=1-2, to=1-4]
\end{tikzcd}
\]

Then the 1st red line is the composition $(h\circ g)\circ f$ and the 2nd red line is the composition $h\circ
(g\circ f)$. Then by associativity, these two functions are equal
\[
	(h\circ g)\circ f = h\circ (g\circ f)
\]
a diagram which satisfies this property for all paths is called a \emph{commutative diagram}. These shall show
up everywhere in this book, and mathematics more generally. They shall be further explored in
section~\ref{sec:categoriesPrelim}.

\subsection*{Injective, Surjective, and Bijective}



There are two special types of functions that deserve to be pointed out: \emph{injective} functions and
\emph{surjeciton} functions. These shall each represent some partial ``reverse'' of a function. Starting off
with injectivity:

\index{Injective}
\begin{defn}{Injective}{injective}
	Let $f: X \rw Y$ be a function. Then $f$ is called \emph{injective}, an \emph{injection}, or
	\emph{one-to-one} if the following condition holds: for
	all $a,b \in X$
	\[
		f(a) = f(b)\ \Rw\ a = b
	\]
	To emphasize a function is injective, a hooked right arrow ($\hookrightarrow$) may be used like so:
	\[
		f: X \hookrightarrow Y \qquad X \xhookrightarrow{f} Y \qquad X\hookrightarrow Y
	\]
\end{defn}

Intuitively, being injective means we can embed the domain $X$ directly into $Y$. As a consequence of this, we can
always ``recover'' $X$ in the following way

\begin{prop}{Injective And Left Inverse}{injectiveAndLeftInverse}
	let $f$ be an injective function. Then there exists a function $\ell f(X) \rw X$ such that
	\[
		\ell\circ f = \id
	\]
	that is $(\ell\circ f)(x) = \id(x) = x$ is the identity function.
\end{prop}

\begin{Proof}
	Define $\ell: Y \rw X$, $\ell(y) = \ell(f(x)) = x$. Since $f$ is injective, this is a well-defined
	function (if it were not injective, then $\ell(f(x))$ would be ambiguous). Then by definition $(\ell\circ
	f)(x) = \ell(f(x)) = x$, as we sought to show.
\end{Proof}

Notice that the inverse function $\ell$ inverted the function \emph{after} applying $f$. What if there is some
function $r: Y \to X$ such that $f$ reverses it? The following shows when this is the case:

\index{Surjective}
\begin{defn}{Surjective}{surjective}
	Let $f: X \rw Y$ be a function. Then $f$ is \emph{surjective}, a \emph{surjection}, or \emph{onto} if for all $y \in Y$,
	there exists an $x$ such that $f(x) = y$, i.e. $f(X) = Y$. To emphasize a function is surjective,
	a double-headed right arrow ($\twoheadrightarrow$) may be used like so:
	\[
		f: X \twoheadrightarrow Y \qquad X \overset{f}{\twoheadrightarrow} Y \qquad X\twoheadrightarrow Y
	\]
\end{defn}

Intuitively, being surjective means that $Y$ is ``represented'' by the elements of $X$, i.e., every element
$y$ has one or many elements of $x$ associated to it, which is why another word for a function to be
surjective is that it is ``onto''. As a consequence of this, we can pullback any element $y \in Y$ to some
representative $x \in X$ such that $f(x) = y$. We can capture this intuition in the following way:

\begin{prop}{Surjective And Right Inverse}{surjectiveAndRightInverse}
	Let $f: X \rw Y$ be a surjective function. Then there exists a (and in fact more generally many) $r: Y \rw
	X$ such that
	\[
		r \circ f = \id
	\]
	that is $(f\circ r)(x) = \id(x) = x$ is the identity function.
\end{prop}

\begin{Proof}
	Define $r: Y \rw X$ by mapping each $y \in Y$ to some element $x \in X$ such that $f(x) = y$ (such an
	element exists since $f$ is surjective). Then by definition for any $y \in Y$, $(f\circ r)(y) = f(r(y)) =
	f(x) = y$ as we sought to show.
\end{Proof}

Another perhaps more intuitive way of seeing how surjections make every $y$ ``represented'' by $X$ is by seeing that
every surjection defines an equivalence class:

\begin{prop}{Surjective And Equivalence Class}{surjectiveAndEquivalenceClass}
	Let $f: X \rw Y$ be a surjective function. Then we can define an equivalence class $\sim_f$ on $X$ by defining:
	\[
		a \sim_f b \qquad\LRw\qquad f(a) = f(b)
	\]
	If $[x] \in X/\sim_f$ is an equivalence class, such that $f(x) = y$, then $[x]$ is called the \emph{fiber}
	over $y$.
\end{prop}

\begin{Proof}
	\begin{enumerate}
		\item For each $a \in X$, $f(a) = f(a)$ by right-uniqueness, hence $a \sim_f a$
		\item If $a\sim_f b$ for $a,b \in X$, then $f(a) = f(b)$. By symmetry of the equivalence relation $=$,
			$f(b) = f(a)$, and so $b\sim_f a$
		\item Can the reader show that if $a\sim_f b$ and $b\sim_f c$, then $a\sim_f c$?
	\end{enumerate}
\end{Proof}

If $f$ is not surjective, this equivalence relation may still be defined on the image of $f$. If $f$ is
injective, then each equivalence class has size $1$. Many functions that are present in mathematics are both
injective and surjective; to give a simple example:
\begin{equation}\label{eq:3setBijective}
		f: \{1,2,3\} \to \{a,b,c\}, \qquad f(n) =  \begin{cases}
		a & n = 1\\
		b & n = 2\\
		c & n = 3
	\end{cases}
\end{equation}
Sch functions are given a name:

\index{Bijective}
\begin{defn}{Bijective}{bijective}
	Let $f: X \to Y$ be a function. Then $f$ is said to be \emph{bijective} or \emph{isomorphic in} \textbf{Set} if $f$ is injective and
	surjective. If $X$ and $Y$ are bijective, then we may write:
	\[
		X\cong_{\textbf{Set}} Y\qquad X\cong Y \qquad X \xrightarrow{\sim} Y
	\]
\end{defn}

The intuition of bijective functions can combine the intuition behind injective and surjective functions,
however they may equivalently be seen as formally defining a notion of ``size''. As seen in
equation~\ref{eq:3setBijective}, two sets that are bijective have the same size. The issue becomes more subtle
when the domain/codomain are infinite:

\begin{example}{Size Of Infinite Sets}{sizeInfiniteSets}
	Define $2\N = \set{2n}{n \in \N}$ and $f: \N \to 2\N$ by:
	\[
		f(n)  =2n
	\]
	Show that $f$ is bijective.
\end{example}

This gives rise to the following definition:
\index{Cardinality}
\begin{defn}{Cardinality}{cardinality}
	Let $X,Y$ be sets. Then $X$ is said to have the same \emph{cardinality} as $Y$ if there exists a bijective
	function $f: X \to Y$. If $X,Y$ have the same cardinality, they are said to have the same size.
\end{defn}

Another important property of bijective functions is that they give rise to a \emph{unique} inverse:

\begin{prop}{Bijectivity And Unique Inverse}{bijectivityAndUniqueInverse}
	Let $f: X \to Y$ be a bijective function. Then there exists a unique function $g$ such that
	\[
		g\circ f = \id_X \qquad f\circ g = \id_Y
	\]

	where $\id_X: X \to X$ and $\id_Y: Y \to Y$ are the identity maps on the respective sets. Often, the
	subscript on the identity maps are dropped and the following abuse of notation is carried out:
	\[
		g\circ f = f\circ g = \id
	\]
	The function $g$ is often denoted $f^{-1}$
\end{prop}

\begin{Proof}
	Since $f$ is surjective, there exists an inverse function $g: Y \to X$ such that $f\circ g = \id_Y$. Since
	$f$ is injective, there can only be one such function (there is only element in the fiber above each $y
	\in Y$). Furthermore, $\im(f) = Y$, hence $g$ satisfies the condition $g\circ f = \id_X$, as we sought to
	show.
\end{Proof}

\begin{titledBox}{Remark: inverse notation abuse}{inverseNotationAbuse}
	The symbol $f^{-1}$ has one other  of notations that the reader should be aware, in particular this
	symbol may represent one of two things depending on the context. Given $f: X \to Y$:
	\begin{enumerate}
		\item if $f$ is bijective, $f^{-1}$ represents the unique inverse function
		\item $f^{-1}: P(Y) \to P(X)$ is the function
			\[
				f^{-1}(T) = S \in P(X) \qquad S = \set{x \in X}{f(x) = y}
			\]
	\end{enumerate}
	the images of $f^{-1}$ are called the \emph{pre-images} of $f$. This notation is particularly prevalent in
	topology where this the pre-image function takes center stage through continuous functions, however it
	will also have many applications in algebra.
\end{titledBox}

Given any function $f: X \to Y$, there are certain natural functions that we may define that allows us to
decompose $f$ as a composition of injective and surjective functions with an intermediary bijective function

\begin{prop}{Function Decomposition}{functionDecomp}
	Let $f: X \to Y$ be a function. Then there exists functions such that
	\[
\begin{tikzcd}
X \arrow[r, two heads] \arrow[rrr, "f", bend left] & X/\sim_f \arrow[r, "\sim"] & \im(f) \arrow[r, hook] & Y
\end{tikzcd}
	\]
	where the 1st function is surjective, the 2nd bijective, and the 3rd injective
\end{prop}

\begin{Proof}
	Define $\pi_f: X \to X/\sim_f$ by
	\[
		\pi_f(x) = [x]_f
	\]
	that is, $\pi_f$ maps $x$ to it's equivalence class. Next, define $F: X/\sim_f \to \im(f)$ by
	\[
		F([x]) = y
	\]
	where $y\in Y$ is chosen so that $f(x) = y$. Finally, define $\iota: \im(f) \to Y$ by
	\[
		\iota(y) = y
	\]
	that is, $\iota$ maps $y$ to itself (since $\im(f) \sse Y$). The reader should verify that these maps have
	the desired properties.
\end{Proof}

This proposition gives an intriguing relation between equivalence relations and functions, namely any function
may be ``represented'' as an equivalence relation, and equivalently any equivalence relation my be
``represented'' by a function.

Using functions, we may define the notion of a multiset:
\index{Multiset}
\begin{defn}{Multiset}{multiset}
	Let $X$ be a set. Then a \emph{multiset} given by elements of $X$ is a function:
	\[
		m: X \to \N
	\]
	which assigns each element $x \in X$ a multiplicity.
\end{defn}

\begin{example}{Multiset}{multiset}
	Let $X = \{a,b,c\}$ and:
	\[
		m(a) = 3 \qquad m(b) = 1 \qquad m(c)  = 0
	\]
	Then $m$ defines a multiset. Commonly, this set is written as:
	\[
		\{a,a,a,b\}
	\]
	As practice, define another multiset $n$ and show how to define the notion of inclusion, union,
	intersection, and so forth.
\end{example}




As a closing remark, it is a common convention in the study of Analysis (one of the major studies of mathematics) that
if we have a function $f: X \rw Y$, we might abuse notation and use a subset of $X$ as the actual domain. For example,
$f: \R \rw \R$, $f(x) = \frac{1}{x}$ is not a function as we've defined the concept, since $f(0) = \frac{1}{0}$ is not
defined \footnote{This would be covered rigorously in a calculus or analysis class}.
If we were to re-write this to be a function, we'd write $f: \R- \{0\} \rw \R$. However, it is often implied in
analysis that we will restrict the domain to the points for which $f$ is defined.

\begin{Exercise}[label=ex:functionsExercises]
\Question Let $f: X \rw Y$ and take $A \sse X$, $B \sse Y$. Show that $f^{-1}(A) \sse B\ \LRw\ A \sse f(B)$

\Question Let $f: X \rw Y$. Then is it true that $f(A\cup B) = f(A) \cup f(B)$? What about $f(A\cap B) = f(A) \cap f(B)$? What if
there were more unions or intersections?

\Question Let $f: X \rw Y$, and define $f^{-1}: P(Y) \rw P(X), f^{-1}(Y) = \set{x \in X}{f(x) \in Y}$. Is $f^{-1}(A\cup B)
= f^{-1}(A)\cup f^{-1}(B)$? What about $f^{-1}(A\cap B) = f^{-1}(A \cap B)$? What if there were more unions or
intersections?

(In a way, equivalence relations capture some idea of ``projecting''. If you have some set $S$ an you want to put
together all equivalent objects, where equivalence here can be thought as ``similar enough so that we can define the
concept of equality''

\Question (partitioning property. Like creating the ``new objects'' where we ignore the ``details'' that is ``consumed'' by the
equivalence relation)

\Question (natural projection?? Give $\pi: G \rw \left\{ 0 \right\}$ a natural projection that always exists?)
\end{Exercise}

\begin{Answer}[ref={ex:functionsExercises}]
  \Question
\end{Answer}


\subsection{Adding Structures to Sets}

So far, we have been extensively working with sets and functions. Sets by themselves have no notion of size, order,
arithmetic operations, geometry, among others possible properties. Many interesting subjects in mathematics deal with
these additional concepts. For this reason, mathematicians often ``augment'', so to speak, sets to have some extra
structure:
\begin{example}{Extra Structure}{extraStructures}
	\begin{enumerate}
		\item (Order Structure). A set $(A, \le)$ with an order $\le$ is called an \emph{ordered set}. Depending on the
			type of order given (ex. partial oder, total order, well order etc.), the set would be called the appropraited
			order (ex. paritally ordere set, totally ordered set, well ordered set, etc.)
		\item (Algebraic structure) The set $\Z$ on its own is only numbers (and so is $\N, \Q, \R$). If we add
			a function $+: \Z \times \Z \rw \Z$, where $+$ acts exactly like the familiar addition we have learnt
			earlier in life, then $(\Z, +)$ now also has addition! We can add more functions to sets, for example $(\Z,
			+, \cdot)$ where $\cdot: \Z\times \Z
			\rw \Z$ does the usual multiplication. If we augment a set with function(s), we call this an \emph{algebraic structure}. This
			whole book will be dedicated to the study of algebraic structures.
		\item (Topological Structures). Another common augmentation of a set is something called a topology. Given a set
			$A$, a set
			$\mathcal{T}_A$ is a set of subsets of $A$ that follows three rules:
			\begin{enumerate}
				\item $\emptyset \in \mathcal{T}_A$ and $A \in \mathcal{T}_A$
				\item if  $\{U_i\}_{i \in I} \sse \mathcal{T}_A$ is some arbitrary subset of $\mathcal{T}_A$ indexed by
					$I$, then $\bigcup_{i \in I} U_i \in \mathcal{T}_A$ (i.e. an
					arbitrary union of sets in $\mathcal{T}_A$ is in $\mathcal{T}_A$)
				\item if $U_1, U_2, ..., U_n \in \mathcal{T}_A$ is a finite subset of $\mathcal{T}_A$, then $\bigcap_i U_i \in \mathcal{T}_A$ (i.e. a finite
					intersections of sets in $\mathcal{T}_A$ is in $\mathcal{T}_A$)
			\end{enumerate}
			Then $(A, \mathcal{T}_A)$ is called a \emph{topology}. If a set is augmented with a sets of sets that follow
			these rules, it's called a \emph{topological structure}. Topology is the natural space in which much of
			Real and Complex Analysis, as well as [non-classical] geometry takes place. It is the ``coarsest''
			structure we add to sets to give sets some notion of locality.
		\item (Measure Structures). Similar to a topological structure, a measure structure is one where
			\begin{enumerate}
				\item $\emptyset \in \mathcal{T}_A$ and $A \in \mathcal{T}_A$
				\item if $U_1, U_2, ..., U_n, ... \in \mathcal{T}_A$, then $\bigcup_i U_i \in \mathcal{T}_A$ (i.e. an
					countable union of sets in $\mathcal{T}_A$ is in $\mathcal{T}_A$)
				\item if $U_1, U_2, ..., U_n, ... \in \mathcal{T}_A$, then $\bigcap_i U_i \in \mathcal{T}_A$ (i.e.
					a countable
					intersections of sets in $\mathcal{T}_A$ is in $\mathcal{T}_A$)
			\end{enumerate}
			This structure is central to generalizing the notion of ``size'' for general sets and is the backbone
			defining the integral\footnote{In particular, the Lebesgue integral}.
		\item (Differential structure) If you have done calculus, then you are familiar with the notion of being
			differentiable. There is a field of study in which we can call ``sets'' differentiable. Given a set $S$,
			there is a collection of functions called an \emph{atlas} (say, $\CA$ represents this collection) which
			gives a set $S$ a ``differentiable structure'', and is usually denoted $(S, \CA)$. The details of what the atlas
			$\CA$ contains is left for a course in differential geometry or diffeology.
		\item *(Scheme structure) This will be a very important structure on a topological space that will
			dominate the readers who will study algebraic geometry. A topological space represents some idea
			of locality on a set. A \emph{sheaf} represents having data on each (open) set in a topological
			space in such a way that we can use local data to glue together global data (think of solving
			a problem by breaking it down into smaller problems and knowing you can glue it back together).
		\item *(Model) In section~\ref{sec:propositions}, we have shown how to make logical statements given sets. The
			full study how sets (usually augmented with other structures) may have statements which are true or provable
			is the domain of \emph{model theory}. Usually, a set $M$ along with some collection of functions and
			relations is augmented in such a way to be able to study the properties of logical statements. This would be
			covered in a class on first order logic or model theory more generally.
	\end{enumerate}
\end{example}

This textbook is dedicated to the study of algebraic structures. It will sometimes be the case that we encounter ordered
structures \footnote{see section~\ref{sec:subGrpLattice}}, and only near the end of the book do we encounter topological
structures (p.~\pageref{df:zariskiTop}) where it is assumed the reader has taken a course covering point-set
topology\footnote{Differentiable structure's will be mentioned to motivate the concept of schemes which are based off
of manifolds}. The reason the eventual mixing of structure is because different domains of mathematics start motivating
each other in further exploration. For example, many geometers and algebraists found parallels between their fields,
and hence algebraic geometry was born. Similarly, it was found that algebra can be used to solve topological problems,
and hence algebraic topology was found. (A word on geometric group theory and something like LCH groups or Amenable
groups?)

\subsection{*Cardinality}

(Here in there in he book, cardinality will come up. For ex, in proof that an injective map from $X \rw X$ where $X$
is finite is surjective, or when we figure out the algebraic numbers are countable, and other results later in the
book)

(classical diagonalization proof?)

(as exercise, do $\{0,1\}^\N$ is bijective to $[0,1]$ by thinking of binary representation?)
\newpage
\section{Number Theory}\label{sec:numbTheory}

Numbers are at the heart of mathematics and algebra. In this section, the basic properties of numbers will be
explored. Some history of these concepts will be given as well to give a sense of why we call these special sets numbers.

Numbers started with our intuitive ability to understand discrete quantities (ex. $5$ rocks, $2$ people). In
the past, it was contentious whether $0$ was a number (the idea being that numbers are supposed to represent
a real quantity), however it is now accepted that $0$ represents the notion of ``nothingness'', ``emptyness'',
or ``inaction''. Ratios of numbers (e.x $3/4$) became important when it was important to be precise on what
part of a whole object we want (ex. half the money). Negative numbers became important when considering other
peoples debts (you owe $2$ sheep, you ``have'' $-2$ sheep). In the times of at least the Greeks, it was
discovered that there are certain numbers that are irrational (ex. $\sqrt{2}$ cannot be written as
$\frac{a}{b}$,
where $a,b$ are natural numbers); the square root of any prime number shown to be irrational. The square-root of
negative numbers remained undefined for awhile until NAME discovered that they are useful in finding roots of
polynomials of degree 3, and it wans't till Euler that square-roots of negative numbers got the symbol
$i := \sqrt{-1}$. It wasn't till later that it was understood there were many more irrational numbers than
rational numbers; the collection of all rational and irrational numbers form the real numbers. The collection
of all numbers $a+bi$, $a,b \in \R$ is known as the \emph{complex numbers}. For reasons that shall be explored
later in the book\footnote{ref:HERE for details}, any further addition to the number system beyond the
complex numbers start to loose properties of numbers, and so we usually consider the following sets as
``numbers'' (ex. $ab \ne ba$ or $(ab)c \ne a(bc)$):
\[
	\N \sse \Z \sse \Q \sse \R \sse \C
\]
We shall heavily explore the properties of all of these number systems, with the notable exception of $\R$.
This is because the real numbers are in many ways the beginning of the exploration of ``taming infinities''
and start moving beyond algebra. In fact, we shall even show in section~\ref{sec:lowenheimSkolem} that it is
impossible to construct $\R$ in an algebraic way.


\subsection{GCD and Euclidean Algoritm}

\index{Divisibility}
\begin{defn}{Divisibility}{divisibilityWithNumbers}
	Let $a,b$ be numbers. Then we say that $a$ \emph{divides} $b$, or $b$ is \emph{divisible by} $a$ if there
	exists a number $k$ such that
	\[
		ak = b
	\]
\end{defn}

Note that divisibility depends on the set of numbers you are working, In $\N$, $3$ is not divisible by $5$,
however in $\Q$, $3$ is divisible by $5$, namely let $k = 3/5$.

\begin{example}{Divisibility Practice}{divisibilityPractice}
	\begin{enumerate}
		\item Show that if $x \mid n_1$ and $ x\mid n_2$, then $x \mid an_1 + bn_2$ for any $a,b \in \Z$.
		\item Show that if $a \mid b$ and $b \mid c$ then $a \mid c$
	\end{enumerate}
\end{example}

From now on in this section, we shall limit ourselves to divisibility in $\Z$\footnote{In
chapter~\ref{cha:eucPIDandUDF}, we shall re-define divisibility in a more general context}.

One special result about integers is that every number has a prime factorizaiton:

\begin{prop}{Prime Factorization Of Natural Numbers}{primeFactorizationNatNum}
	Let $n \in \N$. Then there exists distinct prime numbers $p_1, p_2, ..., p_k$ and natural numbers $n_1,
	n_2, ..., n_k$ such that $n$ can be uniquely broken down into:
\[
	n = p_{1}^{n_1}p_{2}^{n_2}\cdots p_{k}^{n_k}
\]
\end{prop}

\begin{Proof}
	exercise (hint: use induction)
\end{Proof}

This definition can be extended to $\Z$, except we loose uniqueness, the decomposition is unique up to
a multiple of $-1$. This isn't too much of a concern, since $-1$ has the special property of having
a multiplicative inverse, that is we may multiply $-1$ by a number (in this case, itself) to get to $1$:
\[
	-1\cdot -1 = 1
\]
no other integer (besides trivially $1$) can be multiplied by another integer to get to $1$. Since all numbers
can be broken don into multiples of primes, we can define a ay to compare ho similar to numbers are:

\index{GCD!Numbers}
\begin{defn}{Greatest Common Divisor}{greatestCommonDivisor}
	Let $a,b \in \Z$ be to non-zero integers. Then the \emph{greatest common divisor} between these to numbers is the
	greatest number $n$ such that $n$ divides $a$ and $n$ divides $b$.

	The common notation for this fact is that:
	\[
		n \mid a \qquad n \mid b
	\]
\end{defn}

Similarly, we have the dual concept:

\index{LCM!Numbers}
\begin{defn}{Least Common Multiple}{leastCommonMultiple}
	Let $a, b \in \Z$ be non-zero integers. Then the \emph{least common multiple} is the smallest integer $m$ such that
	\[
		a \mid m \qquad b \mid m
	\]
\end{defn}

\begin{example}{Greatest Common Divisor}{greatestCommonDivisor}
	\begin{enumerate}
		\item $\gcd(10,5) = 5$
		\item $\gcd(100, 90) = 10$
		\item $\gcd(45, 49) = 1$
		\item If $p,q$ are distinct prime numbers, then $\gcd(p,q) = 0$
	\end{enumerate}
\end{example}

It is natural that two prime numbers have $1$ as their greatest common divisor for they have no factor in
common. However, it is not necessary for numbers to be prime for their $\gcd$ to be $1$. We give a name to
numbers with no common factors:

\index{Coprime Numbers}
\begin{defn}{Coprime Numbers}{coprimeNumbers}
	Let $a,b \in \N_{>0}$. Then if $\gcd(a,b) = 1$, $a$ and $b$ are said to be \emph{coprime}.
\end{defn}


We may find the greatest common multiple by finding the prime factorization, however this is in general very
inefficient on large numbers. These values can be found in an algorithmic manner relying on the fact that
numbers have a notion of size\footnote{We shall see in chapter~\ref{cha:eucPIDandUDF} that not all systems in
which we can define a GCD will allow for a euclidean algorithm}

\index{Euclidean Algorithm}
\begin{thm}{Euclidean Algorithm}{euclidAlg}
	Given two natural numbers $a,b$ we can find their greatest common divisor by the following technic:
	\begin{enumerate}
		\item Let $a > b$. Put them into the equation $a = bk_1 + r_1$, where $k$ is the amount of times $b$ fits into $a$,
			and $r_1$ is the remainder.
		\item Once you find these values, all but $k$ shift over to the left, and you repeat the process to find $r_2$:
			\begin{center}
				\begin{align*}
					b &= r_1k_2 + r_2\\
					r_1 &= r_2k_3 + r_3\\
					r_2 &= r_3k_4 + r_4
				\end{align*}
			\end{center}
		\item The algorithm continues until $r_n = 0$ (which it will reach since $a,b$ where finite). The final non-zero
			remainder $r_{n-1}$ is the greatest common divisor.
	\end{enumerate}
\end{thm}

\begin{Proof}
	exercise
\end{Proof}

As an example:

\begin{example}{Euclidean Algorithm}{eucAlgEx}
	Let $a = 20$, $b = 13$. Then:
	\begin{align*}
		20 &= 1\cdot 13 + 7\\
		13 &= 1\cdot 7 + 6\\
		7 &= 1\cdot 6 + 1\\
		6 &= 6\cdot 6 + 0\\
	\end{align*}
	so $\gcd(20,13) = 1$, (the number above $0$ in the last step). To find the $ax + by = 1$ formula, we reverse the
	process by converting the equation above in the following form:
	\[
		7 - 6 = 1
	\]
	so, put the remainder on the right hand side. Then, isolate the next equation in in the appropriate ways to
	substitute into this one: $7 = 13-6$, $6 = 13-7$. Since $6 =13-7$, $7 = 13-(13-7) = 7$, so
	\[
		7 - (13-7) = 1
	\]
	then, we repeat the same process with the step above: $7 = 20-13$, so:
	\[
		20-13-(13-(20-13)) = 1
	\]
	implying, we get
	\[
		2\cdot 20 - 3\cdot 13 =1  \qquad (40-39 =1)
	\]
	giving us our desired equation.
\end{example}

Using the euclidean algorithms, we may always find the solution to the following interesting linear equation:

\begin{thm}{Extended Euclidean Algorithm (B\'ezout's Identity)}{extendedEuclidAlg}
	Let $a,b \in \Z$. Then there exists integers $x,y \in \Z$ such that
\[
	ax +by = \gcd(a,b)
\]
\end{thm}

\begin{Proof}
	exercise (reverse the process)
\end{Proof}

\begin{example}{B\'ezout's Identity}{b'ezoutIdentity}
	Recall that $\gcd(20, 13) = 1$. To find the $ax + by = 1$ formula, we reverse the
	process by converting the equation above in the following form:
	\[
		7 - 6 = 1
	\]
	so, put the remainder on the right hand side. Then, isolate the next equation in in the appropriate ways to
	substitute into this one: $7 = 13-6$, $6 = 13-7$. Since $6 =13-7$, $7 = 13-(13-7) = 7$, so
	\[
		7 - (13-7) = 1
	\]
	then, we repeat the same process with the step above: $7 = 20-13$, so:
	\[
		20-13-(13-(20-13)) = 1
	\]
	implying, we get
	\[
		2\cdot 20 - 3\cdot 13 =1  \qquad (40-39 =1)
	\]
	giving us our desired equation.

	Try doing the same with $220$ and $38$.
\end{example}


B\'ezout’s Identity has the following neat interpretation: primes are fundamentally a multiplicative decomposition of
numbers: any $n$ is equal to $p_1^{n_1}\cdots p_{k}^{n_k}$ for appropriate primes $p_i$ and powers $n_i$. However, they
do have an additive relation between them, the B\'ezout identity. If we have two primes, we can (up to appropriate
constants) add them to get to the multiplicative identity $1$. In this way, we have a property of primes given via
addition.

(should I add other number theory identities? Are there any other useful ones for this book? Perhaps just as exercises?
They might be useful way later in algebraic number theory)

\begin{Exercise}[label=ex:gcdAndLcm]
  \Question Let $p,q,r \in \N$ be prime numbers greater than $3$. Show that $3$ divides $p^2+q^2+r^2$
\end{Exercise}
\begin{Answer}[ref={ex:gcdAndLcm}]
  \Question
\end{Answer}
\subsection{Modular Arithmetic’s}\label{sec:modArithmetics}

When counting the hours on a clock-face, we start with $12$ o-clock, then add an hour to make it $1$ o-clock,
then another hour to make it $2$ o-clock, until we circle back to 12 o-clock again, at which point we start
counting again to $1$ o-clock. For the rest of this section, we will say that $12$ o-clock is $0$ o-clock for
reasons that would be explained soon. Notice how the arithmetic's of the clock different from the usual number
system where for each $x \in \N$, there exists an $n$ such that $n > x$. Hours on a clock could add to be
zero, for example:
\[
	7+5=  0
\]
In this section, we rigorously define this notion

\index{Modular Arithmetics}
\begin{defn}{Modular Arithmetic's}{modularArithmetics}
	Let $n \in \N_{>0}$. Define the equivalent relation on $\Z$:
	\[
		a\sim_n b \quad \LRw \quad a-b = kn\text{ for some $k \in \Z$}
	\]
	The collection of these equivalences classes is denoted $\Z/n\Z$\footnote{This notation shall be
	elucidated on in chapter~\ref{cha:quotientGrps}} and is called the \emph{integers modulo n}. The number $n$
	is called the \emph{modulus} and if $a\sim_n b$, we say that \emph{a is congruent to b} modulo $n$. If
	$a\sim_n b$, then it is denoted
	\[
		a\equiv b \mod n \qquad \text{or} \qquad a \equiv_n b
	\]
	The equivalence classes may be denoted as:
	\[
		[a],\ [a]_n,\ \oln{a},\ \oln{a}_n
	\]
	where the $n$ is dropped if the context permits it.
\end{defn}

Intuitively, this equivalence relation says that the difference of two numbers is a multiple of $n$. When $n
= 12$, then this is the usual clock arithmetic's. When $n$ is any other number, the reader may imagine
a clock-face with $n$ digits.

Naturally, it has to be proven that $\sim_n$ is indeed an equivalence relation:

\begin{prop}{Modular Arithmetic's Well-defined}{modArithmeticsWelldefined}
	Let $\sim_n$ be defined as in \cref{df:modularArithmetics}. Then $\sim_n$ is indeed an
	equivalence relation
\end{prop}

\begin{Proof}
	Let $x,y,z \in \Z$ and let $n \in \N_{>0}$ be any positive natural number.
	\begin{enumerate}
		\item $x \equiv x \mod n$ since $x-x = 0 = 0\cdot n$
		\item If $x\equiv y\mod n$ then there exists a $k$ st $x-y = kn$. Then $y-x = -kn = (-k)n$. Then by
			definition $y\equiv x \mod n$
		\item If $x\equiv y \mod n$ and $y\equiv z \mod n$, then $x-y = kn$ and $y-z = \ell n$ for some
			$k,\ell \in \Z$. Adding these two equations, we get:
			\[
				x-z = x-y+y-z = (k+\ell)n
			\]
			and so $x\equiv z \mod n$
	\end{enumerate}
	Since all the conditions of an equivalence relation is satisfied, $\sim_n$ is indeed an equivalence
	relation, as we sought to show.
\end{Proof}

\begin{example}{Modular Arithmetic's}{modularArithmetics}
	Verify the following:
	\begin{enumerate}
		\item $35 \equiv 11 \mod 12$
		\item $2\equiv -3 \mod 5$
		\item If $a\equiv b\mod n$, then $a+t \equiv b+t \mod n$ for any $t \in \Z$
		\item If $a\equiv b\mod n$, then $a^t \equiv b^t \mod n$ for any $t \in \N$
	\end{enumerate}
\end{example}

\begin{prop}{Properties Of Modular Arithmetic's}{propModularArithmetics}
	Let $n \in \N_{>0}$ and $a_1\equiv a_2 \mod n$, $b_1\equiv b_2\mod n$. Then:
	\begin{enumerate}
		\item $a_1+b_1\equiv a_2+b_2\mod n$
		\item $a_1-b_1\equiv a_2-b_2\mod n$
		\item $a_1b_1 \equiv a_2b_2\mod n$
	\end{enumerate}
	Furthermore:
	\begin{enumerate}
		\item $a_1+b_1\mod n = a_1\mod n + b_1\mod n$ (in different notation, $\oln{a_1+b_1} = \oln{a_1}
			+ \oln{b_1}$)
		\item $a_1b_1\mod n = (a_1\mod n)(b_1\mod n)$ (in different notation, $\oln{a_1b_1}
			= \oln{a_1}\oln{b_2}$)
	\end{enumerate}
\end{prop}

\begin{Proof}
	This is a good exercise
\end{Proof}


\begin{prop}{Cancellation Properties Of Modular Arithmetic's}{cancellationPropModularArithmetics}
	Let $n \in \N_{>0}$. Then:
	\begin{enumerate}
		\item if $a+k\equiv b+k\mod n$ for any integer $k \in \Z$, then $a\equiv b \mod n$
		\item if $ka\equiv kb\mod n$ for $k$ coprime with $n$ ($\gcd(k,n) = 1$) then $a\equiv b \mod n$
		\item if $ka\equiv kb \mod kn$ for some $\N_{>0}$, then $a\equiv b \mod n$
	\end{enumerate}
\end{prop}

\begin{Proof}
	good exercises
\end{Proof}

\subsection*{Important Results}

The following are important properties of the integers modulo $n$:

\index{Fermat's Little Theorem}
\begin{thm}{Fermat's Little Theorem}{fermatLittleThm}
	Let $p$ be prime and $a \in \N_{>0}$ such that $p\nmid a$. Then
	\[
		a^{p} \equiv a \mod p
	\]
\end{thm}

\begin{Proof}
	First, we can assume that $0 \le a \le p-1$, for we may always reduce $a$ due to the properties of modular
	arithmetic's (hint: use the multiplicative properties of modular arithmetic's). Next, we may reduce this
	assertion to
	\[
		a^{p-1}\equiv 1 \mod p
	\]
	For, if the above hold, then multiplying both sides by $a$ gives the desired result.

	Start by list out the first $p-1$ positive multiples of $a$:
	\begin{equation}\label{eq:aTermListing}
				a, 2a, 3a, ..., (p-1)a
	\end{equation}
	I first claim that no two of these are equal $\mod p$. For the sake of contradiction, let's say two were
	so that $ra \equiv sa\mod p$ for some $1 \le r,s \le p-1$ where $r\ne s$. Then since $p$ is prime and $a
	> 0$, $\gcd(a,p) = 1$ an hence by \cref{pr:propModularArithmetics}:
	\[
		r\equiv s \mod p
	\]
	Since $0 \le r,s \le p-1$, this implies $r = s$; contradiction our earlier claim. Hence, each term in
	list~(\ref{eq:aTermListing}) is congruent (modulo $p$) to a unique number $1, ..., p-1$.

	Now, multiply the terms in list~(\ref{eq:aTermListing}). Then by
	\cref{pr:propModularArithmetics}:
	\[
		(1a)(2a)\cdots((p-1)a) \equiv (1)(2)\cdots(p-1)\mod p
	\]
	which, when simmplifying by combining terms, is:
	\[
		a^{p-1}(p-1)! \equiv (p-1)!\mod p
	\]
	Since $p$ is prime, each number $1 \le k \le p-1$ is coprime to $p$, and hence by
	repeated application of \cref{pr:cancellationPropModularArithmetics}:
	\[
		a^{p-1} \equiv 1 \mod p
	\]
	but now, we may simply multiply both sides by $a$ to get $a^{p} \equiv a \mod p$, as we sought to show.
\end{Proof}

The reader may ask if this can be generalized in some some way to modulo $n$. The key would be to generlize
the step in which each number $1 \le k \le p-1$ is coprime with $n$, and indeed if $\phi(n)$ represents the
number of numbers that are coprime with $n$, then the identity can be generalized to:
\[
	a^{\phi(n)} \equiv 1 \mod n
\]
This proof is best in the context of group theory, and so shall be left alone till then, however the
persistent reader has all they need to attempt to directly prove this result.

(There are more cool proofs, but many rely on group theory or field theory. I'd like to incldue them in
certain sections of the book: \href{https://en.wikipedia.org/wiki/Modular_arithmetic}{here})


\begin{Exercise}[label=ex:numberTheory]
  \Question Let $a\equiv b \mod n$. If $k\mid a,b$, is it the case that $(a/k) \equiv (b/k)\mod n$?
\end{Exercise}
\begin{Answer}[ref={ex:numberTheory}]
  \Question No, take $a = 6$, $b = 0$, and $k = 3$.
\end{Answer}

\subsection{Counting Problems}

(Also introduce the binomial theorem?)
(somehow introduce the binomial theorem. Something else too?)
\newpage
\section{Graphs}\label{sec:graphs}

mention how this is section is not necessary to read for groups, but it will come back up when working with particular
perspective on groups. It can be revisited once you make it there (perhaps make the graph intuition of groups a series
of exercises?)

(See this link: (commented for compiling purposes) %https://tjyusun.com/mat202/sec-graph-isomorphisms.html
)

We first introduce a more concrete definition of a graph in order to be able to work with it unambiguously

\index{Graph}
\begin{defn}{Graph}{graph}
	Let $V = \{ v_1, v_2, ..., v_n, ...\}$ be a set of elements representing the \emph{vertices} (also called \emph{nodes}
	or \emph{points}). Let
	\[
		E = \set{\{v_i, v_j\}}{v_i,v_j \in V}
	\]
	be the set representing the \emph{edges} of the graph (where there need not be a set $\{v_i, v_j\}$ for every pair of vertices in
	$V$). Then $\mathfrak{G} = (V, E)$ is called a \emph{graph}.

	If the edges are \emph{pairs} instead of sets, then we call $\mathfrak{G}$ a \emph{direct graph}
\end{defn}

\begin{example}{Graphs}{graphsEx}
	put different types of graphs as example? infinite, connected, cyclic, tree, completed
	\begin{enumerate}
		\item here + visual
		\item here + visual
		\item here + visual
		\item here + visual
	\end{enumerate}
\end{example}

(define endpoints of a vertex, adjacent vertices, and if $e = \{u,v\}$ $e$ is also called the incident of $u$ and $v$)

(define simple graph: $E$ has no repeated elements. If $E$ can have repeated elements, it's a multigraph)

%\begin{example}{graph exercises}{graphExercices}
	%\begin{enumerate}
		%\item tree is a-cyclic
		%\item probably counting vertices and edges problems?
		%\item other simple and fun graph theory problems?
	%\end{enumerate}
%\end{example}

\index{Graph Isomorphism}
\begin{defn}{Graph Isomorphism}{graphIso}
	Let $\mathfrak{G} = (V_{\mathfrak{G}}, E_{\mathfrak{G}})$ and $\mathfrak{H} = (V_{\mathfrak{H}}, E_{\mathfrak{H}})$
	be two graphs. A function $f: V_\mathfrak{G} \rw V_\mathfrak{H}$ is called a \emph{graph isomorphism} if for every
	vertex adjacent vertices $v_i, v_j \in V_{\mathfrak{G}}$, we have
	\[
		(v_i, v_j) \in E_{\mathfrak{G}}\ \LRw\ (f(v_i), f(v_j)) \in E_{\mathfrak{H}}
	\]
	that is, $f$ preserves adjacencies between vertices.

	If $\mathfrak{G}$ and $\mathfrak{H}$ are isomorphic, we write $\mathfrak{G} \cong \mathfrak{H}$, and we say that
	$\mathfrak{G}$ is isomorphic to $\mathfrak{H}$
\end{defn}

\begin{example}{examples of Graph Isomorphisms}{graphIsoEx}
	here
\end{example}

These will come back in a exercise ref:HERE
\subsection{Lattices}\label{sec:lattice}

(These I think are important because they will be periodically mentioned. From the beginning with the partial ordering
of subsets via inclusion, to using Zorn's lemma, to group/ring/module lattices, to the correspondence between the latice of divisibility in UFDs and
multi-sets of irreducible elements, to the many galois correspondences we shall see, to Quivers)


Grillet has a good section on Lattices, so does Lang.
%(should these be included somewhere?  \href{https://en.wikipedia.org/wiki/Lattice_(order)}{lattices} Or perhaps under graphs? )


\begin{Exercise}[label=graphExercise]
	\Question tree is a-cyclic
	\Question probably counting vertices and edges problems?
	\Question other simple and fun graph theory problems?
\end{Exercise}
\begin{Answer}[ref={graphExercise}]
	\Question Do this that and the other
	\Question This formula
	\Question No answer with ambiguous question \smiley
\end{Answer}
\newpage
\section{Categories}\label{sec:categoriesPrelim}

% (maybe at the end give the definition of a category just to expose the reader to it?)

In the last couple of section, we have been focusing on understanding mathematical objects, in particular, we
tried constructing different types of objects (like sets and functions, numbers, graphs, and sets with binary
relations) and tried showing some properties of these objects. In this section, we will focus less on the
mathematical object, and more on the \emph{relation} between mathematical objects that preserve some of the
same ``structure''. The nouns ``structure'' and ``relation'' will be given more precise meaning through the
use of \emph{categories}.

Categories have a reputation to be difficult to understand -- a reputation I (the author of this book) feel
comes from an overly abstract introduction at too early a stage in the pedagogical journey. The hardest part
about categories is a shift in perspective in how to think about mathematical objects, and having the sample
space to intuit why this shift brings great value. Thus, before getting into the theory of category theory,
let's first give an analogy to get a more intuitive understanding in what we are trying to achieve:

\begin{quote}
	Let's take any Language (ex. English). We can study any given word to get some properties of it (for example, its
	grammatical properties, syntactic correctness, its phonetic properties in a dialect, and so forth). These
	give information about words, however what
	we usually care about is not the properties of a word, but how words \emph{relate} to other words. The meaning
	behind a sentence doesn't come just from the individual words, but from how the words ``come together'': This is
	called semantics~\footnote{Semantics is the study of the meaning that sentences (or any discourse) carries}. As an
	example of semantics, we can
	even see that the meaning of words is defined in relation to other words in a sentence; for example the
	word \emph{cold} usually means the absence of heat, but in the sentence ``those people are giving
	him a cold stare'', I am not referencing to the temperature of their gaze, but some lack of friendliness in
	their feelings towards him.

	If we look at language through the lens of the relation between words, we can get some insights on how language is
	used. An excellent example of this is translation. Let's say we're trying to translate from English to Sanskrit (or
	really any two languages). Both English and Sanskrit are languages, and so they have words (i.e. the objects) and
	rules (i.e. the relation between objects). Don't worry for a moment about exceptions to rule in particular languages
	(If it helps the reader, pretend they don't exist for the sake of the analogy), the main idea is that English and Sanskrit (and all
	languages ) have words and rules that govern their use. When we compare English to Sanskrit, what we will do is see
	how the relations of the words in English (ex. conjugations, articles, etc.) link to those in Sanskrit (ex.
	declension, compounds, etc.). In a sense, we don't care about what the particular words of the respective
	language are. What we care is to find the word form one language that will create the same relations in the other
	language.

	For example, If I say ``dog'' in English, then the sentences ``The dog barked'', ``I took the dog for a walk'', or
	``Many dogs like to chase cats'' all make sense, while ``I dog the took for a walk'' is not a sentence that makes
	sense, since the position of the word ``dog'' and ``take'' are wrong given the grammatical rules in the English
	languages. In Sanskrit, the grammatical rules
	are different: If we directly translate the sentence ``I dog the took for a walk'' to Sanskrit, the
	sentence will actually make sense! That is because the rules of Sanskrit are different than English. In fact, it is
	a little more subtle than this: There is no ``a'' (like ``a walk'')  or ``the'' in Sanskrit -- if you're translating from
	English to Sanskrit, we can't replace ``a'' with a Sanskrit word (since non exist). Instead, we will make sure that
	the words we choose in Sanskrit when translating the English sentences will have the appropriate \emph{relations} in order
	to capture the same meaning
\end{quote}

The entire point of this analogy is to show that it is useful to understand objects not how they ``intrinsically'' are,
but how they ``in relation'' to other objects \footnote{The philosophy behind this idea is called \emph{structuralism}}.
The analogy showed that an advantage of this approach comes from a better understanding of how to translate
sentences. Similarly, thinking about mathematical objects by how they relate to one another will let us be
able to compare object by seeing if one mathematical object has the same type of relations with other objects.

In fact, we have already done this before: when we worked with ordered structures (ex. $(\N, \le)$),  we studied objects by how they
related to other objects. We can even make definitions based on how objects relate. For example, in $(\N, \le)$, define the \emph{smallest}
to be the object $x$ such that for every other element $y \in \N$, $x \le y$. Then it is clear that $0$ is the smallest
object. Hover, in $(\Z, \le)$ no object is the smallest object, because for any $x \in \Z$, there exists $x-1 \in \Z$
such that $x-1 \le x$. Another example of a definition based on how an object relate, let $X$ be the set for which any
function from $X$ is \emph{always} injective. Then $X$ is any singleton (exercise ref:HERE).

In the following section, we will try to more rigorously define a way to talk about the relation between objects

\subsection{Important Concepts}

This section is very much a cursory glance of the elementary concepts of category theory for the sake of
exposure rather than for internalizing the definitions and working through problems. Let's start by giving a
high-level idea of that a category is. This definition should hold both information about objects and the
relation between objects:

\begin{defn}{Pre-Definition Of Category}{pre-DefinitionCategory}
	A Category \textbf{C} is a collection\footnote{the word \emph{collection} is used instead of \emph{set} for a technical
	reason. More on that after this definition} of objects and relations between objects:
	\begin{enumerate}
		\item We denote the collection of objects of \textbf{C} by $\operatorname{Obj}(b)$.
		\item For any to objects $A,B \in \operatorname{Obj}(\bf{C})$, $\Hom_{\bf{C}}(A,B)$ denotes the collection of
			relation between objects, which we call \emph{morphisms}. In particular, if $f \in \Hom_{\bf{C}}(A,B)$,
			then $f$ is called a morphism. Morphism are \emph{directional}, meaning that if $f \in
			\Hom_{\bf{C}}(A,B)$, it is \emph{not} the case that $f \in \Hom_{\bf{C}}(B,A)$ (unless $A = B$)

			If $f \in \Hom_{\bf{C}}(A,B)$, we shall write $f: A \to B$ to emphasize the directionality.
		\item If $f \in \Hom_{\bf{C}}(A,B)$, and $g \in \Hom_{\bf{C}}(B,C)$, then there must exist a morphism $h
			\in\Hom_{\bf{C}}(A,C)$
			that we associate $g$ and $f$; we'll write $h = gf$, $h = g\cdot f$, or $h = g\circ f$, and call
			$h$ the \emph{composite morphism} of $f$ and $g$.
	\end{enumerate}
\end{defn}

The fact that we said \emph{collection} instead of \emph{set} is because of a quirk of logic. In particular,
once we properly define what a category is, we will show that sets and functions between sets will form a
category. Then $\operatorname{Obj}(\textbf{Set})$ will have to be \emph{a set of all sets}, but it can be
proven that no such set exists! This is called \emph{Russel's Paradox}, and is addressed in
chapter~\ref{cha:catApproach}. The morphisms in $\Hom_{\bf{C}}(A,B)$ cannot be chosen arbitrary, they shall have to satisfy
the following properties:
\begin{enumerate}
	\item (identity morphism) For each $A \in \Obj(\bf{C})$, there must at least exists morph $1_A \in
		\Hom_{\bf{C}}(A,A)$ such that for any
		$f \in \Hom_{\bf{C}}(A,X)$ and $g \in \Hom_{\bf{C}}(X,A)$ (for any $X \in \Obj(\bf{C})$, we have:
		\[
			f\id_X = f \qquad \text{and} \qquad \id_X g =g
		\]
	\item (associativity) For $f \in \Hom_{\bf{C}}(A,B)$, $g \in \Hom_{\bf{C}}(B,C)$ and $h \in \Hom_{\bf{C}}(C,D)$,then if $\alpha = gf$ and
		$\beta = hg$, then:
		\[
			h\alpha = \beta f \qquad \text{i.e.} \qquad h(gf) = (hg)f
		\]
\end{enumerate}

The motivation behind these rules will come in chapter~\ref{cha:catGrp}.

\begin{example}{Category of \textbf{Set}}{preSecCatofSet}
	The prototypical examples of a category is \textbf{Set} with object as sets and morphisms as
	functions, namely for any two sets $X,Y \in \Obj(\bf{Set})$, $\Hom_{\bf{Set}}(X,Y)$ is the collection
	of all functions from $X$ to $Y$. For any $X \in \Obj(\bf{Set})$, the identity function is the
	identity morphism, and section~\ref{sec:functions} showed functions are associative.


	We shall cover more examples when there is more material to cover in chapter~\ref{cha:catGrp} and more
	advanced examples in part~\ref{part:catTheory}.
\end{example}

Let us now go over some important concepts that shall come up continuously throughout the entire textbook.

\index{Isomorphism}
\begin{defn}{Isomorphism}{preCatIso}
	Let \textbf{C} be a category (the reader can keep in mind \textbf{Set}). Then we shall say that $X,Y \in
	\Obj(\bf{C})$ are \emph{isomorphic} if there exists morphisms $f \in \Hom_{\bf{C}}(X, Y)$ and $g \in
	\Hom_{\bf{C}}(Y,X)$ such that $fg = \id_X$ and $gf = \id_Y$. We shall sometimes abuse notation and simply
	write $fg = gf = \id$. If $X,Y$ are isomorphic, we shall write:
	\[
		X\cong_\bf{C} Y
	\]
\end{defn}

\begin{example}{Isomorphisms in Set}{isosinSet}
	Let $\bf{C} = \bf{Set}$. Show that $X,Y$ are isomorphic if and only if $|X| = |Y|$. This shows that the relations in
	the relations in the category \textbf{Set} ``capture'' or ``remember'' cardinality information of sets. If
	we had different morphisms, we can have different types of information.

	Recall that what we care about in category theory are the relations between objects and not necessarily
	the objects themselves. As a consequence, results proven in \textbf{Set} shall be true for all sets of the
	same cardinality. More generally, we shall see that results proven in category for an object $X$ shall be
	true \emph{up to isomorphism}, that is if they are true for $X$, they are true for $Y$ such that $X
	\cong_{\bf{C}} Y$.

	The reader may think to themselves that cardinality information is not much to preserve, for example we
	may have two sets of the same cardinality that satisfy different distinct predicates. This is because \textbf{Set} is a very simple category and each morphism holds little
	information\footnote{The knowledgeable
	category theorist may complain that by Yoneda's lemma, any locally small category can embed into
	\textbf{Set} and so the category is far from being simple. The aforementioned
	simplicity is from the basic description and isomorphisms classes rather than the complicated
	subcategories we may produce}. We shall often extend the objects and morphism of \textbf{Set} to hold more
	information.
\end{example}

The next important idea is that of \emph{commutative diagram}. These shall represent how we can relate
objects:

\index{Commutative Diagram}
\begin{defn}{Commutative Diagram}{comDiagramPreCat}
	Let \textbf{C} be a category. Then a \emph{commutative diagram} is a graph where each node is an object $X
	\in \Obj(\bf{C})$ and each edge is a morphism (with appropriate domain and codomain) where for each closed
	loop in the graph (disregarding orientation), the composite morphisms must agree.
\end{defn}

To elucidate the last condition, consider:
\[
\begin{tikzcd}
                                   & Y \arrow[rd, "g"] &   \\
X \arrow[rr, "h"'] \arrow[ru, "f"] &                   & Z
\end{tikzcd}
\]
This forms a closed path, and hence composition must agree, and so $gf = h$. Another common example is:
\[
\begin{tikzcd}
X \arrow[r, "f"] \arrow[d, "h"] & Y \arrow[d, "g"] \\
Z \arrow[r, "i"]                & W
\end{tikzcd}
\]
Then:
\[
	gf = ih
\]
As we care about the relations between objects, commutative diagrams will be the building for most categorical
theorems. A lot of results shall come down to saying ``the following diagram commutes''. Let us apply this to
the first concept of category theory: \emph{universal properties}. A universal property shall often be
a property that captures the \emph{essence} of a predicate in such a way that we can ask whether an object in
different categories have or don't have the property. This is part of where the word \emph{universal} comes
from; these shall represent properties that can be described \emph{universally between categories}. The full
generality of universal properties is something we shall delay for chapter~\ref{cha:catGrp}. For now, let us give an
important example of the \emph{universal property of quotients}. We shall show what they exist \textbf{Set}, and
in each part of the book we shall see how the majority of categories we work with shall have objects
satisfying this universal property.

The name suggests some notion of ``dividing''. This is certainly strongly related to dividing in the
traditional numerical sense, however the connection to dividing would be too large of a build up for the
expository nature of this section, and so it will have to be relegated to chapter~\ref{cha:eucPIDandUDF}.
Instead, the reader should think of \emph{equivalence relations} of a way of dividing up a set. The universal
property of quotients in \textbf{Set} shall show that there is an object that represents the equivalence
relation:

\index{Universal Property!Quotient Sets}
\begin{thm}{Universal Property Of Quotient Sets}{uniPropQuotient}
	Let $A$ be a set and $\sim$ be an equivalence relation on $A$. Let $f: A \rw Z$ be any set function such that $a\sim
	b \Rw f(a) = f(b)$. Then there exists a unique map $\phi$ from $A/\sim$ to $Z$ such that the following diagram commutes
	\[
		\begin{tikzcd}
		A/\sim \arrow[r, "F"]             & Z \\
		A \arrow[u, "\pi"] \arrow[ru, "f"] &
		\end{tikzcd}
	\]
\end{thm}

\begin{Proof}
	Let's say $f: A \rw Z$ respected where $a\sim b \Rw f(a) = f(b)$. we want to ask if there exists a unique map such that
	\[
		f = F\circ \pi
	\]
	that is
	\[
		f(x) = F(\pi(x))
	\]
	Consider the map $F: (A/\sim) \rw Z$, $F([x]) = f(x)$. This map is well-defined since if $[x] = [y]$, then $x \sim y$
	so $f(x) = f(y)$.

	Furthermore, this map is unique. If $F, F': (A/\sim) \twoheadrightarrow Z$ were two well-defined
	maps, then it must be that $F(\pi(x)) = f(x) = F'(\pi(x))$ making $F = F'$, showing uniqueness, completing the
	proof.
\end{Proof}

Essentially, given an equivalence relation $\sim$ on a set $A$ it can be ``represented'' by the object $A/\sim
\in \Obj(\bf{Set})$ which captures the right amount of information such that any other object that tries to
captures it is always linked to $A/\sim$. In other words, $A/\sim$ is in a sense always in an object that
tries to represent the equivalence classes of $A$.

This will conclude our quick exposition of category theory. In chapter~\ref{cha:catGrp}, we shall properly
introduce the theory of category theory at the ``elementary level''. In part~\ref{part:catTheory}, we shall
properly cover category theory in its full generality and introduce the important major results that students
ought to carry throughout their mathematical career.

% \begin{enumerate}
% 	\item quotient objects.
% 	\item products (fibered product as an advanced example)
% 	\item An attempt at intuition behind universal properties? Perhaps the best first example is when we see free
% 		objects?????
% 	\item when talking about the category of set, you can say that since we care about the relation between objects, we
% 		loose some details about sets. We won't really know where elements from one set map to another, just how they
% 		map. So, in this sense, the cardinality of sets becomes really important, since that determines what types of
% 		maps we can make (ex. injective, surjective).
% 		\begin{enumerate}
% 			\item Looking at sets in terms of their cardinality will actually permeate a lot of what we'll do. (Can you
% 				figure out how to explain this here??) Instead of saying $F_S$ for a free-group, we'll write $F_n$ where
% 				$|S| = n$, since every set $S$ with the same cardinality are isomorphic in \textbf{Set} and free groups are
% 				defined on a set \emph{up to isomorphism} so we will write $F_S$. Similarly for $S_n$, $\R^n$, and so
% 				on. Basically, wherever we see a number, we can interpret that as a secret use of this ``up to
% 				isomorphism in \textbf{Set}'' fact.
% 		\end{enumerate}
% \end{enumerate}

\newpage
\section{Matrices}\label{sec:matrices}

This section can be comfortably skipped until chapter~\ref{sec:matrixRepOfLinTran}. It will also be referenced in
section \ref{sec:matrixGrp} for examples of groups, but it is skippable without any loss (since it's just one out of
over a dozen example of groups that will be presented).

In this book, we will work with many structures. Some of them will be very different from one another, while others
share striking similarities. Many of these more similar mathematical structures will be able to be ``represented'' by
matrices. We will define more rigorously for something to be ``representable'' by matrices later -- for now we want to
gain some basic appreciation and understanding of how to use matrices.

(History going back to China, this is the early motivator for matrices compared to the ``new'' one of being a useful
classification theorem. Define a matrix. Show it's used to solve systems of linear equations. Explain Gaussian
elimination)

(Later, Gauss developed a way to multiply matrices. The intuition shall come in chapter~\ref{cha:vecSp})




\newpage
\printbibliography
\end{document}
